% =============================================================================
% HardwareX-style report -- EC2 (complete version)
% Mechatronic Systems Design 2026-2, UTEC
% Uses only packages already present in the project.
% Text in RED marked [TO COMPLETE: ...] must be filled in by the authors.
% =============================================================================

\documentclass[11pt,letterpaper]{article}

% -----------------------------
% Encoding and typography
% -----------------------------
\usepackage[T1]{fontenc}
\usepackage[utf8]{inputenc}
\usepackage{lmodern}
\usepackage{microtype}
\usepackage{graphicx}
\usepackage{float}
\usepackage{comment}
\usepackage{amsmath,amssymb}

% -----------------------------
% Page layout and section style
% -----------------------------
\usepackage[margin=1.5cm]{geometry}
\usepackage{titlesec}
\titleformat{\section}[block]{\bfseries}{\thesection.}{0.5em}{}
\titleformat{\subsection}[block]{\bfseries}{\thesubsection}{0.5em}{}

\setlength{\parindent}{0pt}
\setlength{\parskip}{8pt}

% -----------------------------
% Lists
% -----------------------------
\usepackage{enumitem}
\setlist[itemize]{topsep=3pt,itemsep=2pt,parsep=0pt}
\setlist[enumerate]{topsep=3pt,itemsep=2pt,parsep=0pt}

% -----------------------------
% Tables
% -----------------------------
\usepackage{array}
\usepackage{tabularx}
\usepackage{booktabs}
\usepackage{longtable}

\newcolumntype{Y}{>{\raggedright\arraybackslash}X}
\newcolumntype{L}[1]{>{\raggedright\arraybackslash}p{#1}}
\renewcommand{\arraystretch}{1.18}

% -----------------------------
% Colors and hyperlinks
% -----------------------------
\usepackage{xcolor}
\usepackage{xurl}
\usepackage{hyperref}
\hypersetup{
    colorlinks=true,
    linkcolor=blue,
    citecolor=blue,
    filecolor=magenta,
    urlcolor=blue
}
\urlstyle{same}

\definecolor{mastheadgray}{RGB}{230,230,230}
\definecolor{elsblue}{RGB}{0,114,188}

% Items the authors still have to complete (printed in red)
\newcommand{\TODO}[1]{\textcolor{red}{\textbf{[TO COMPLETE: #1]}}}

% -----------------------------
% Author biography (shows a placeholder if the photo is missing)
% -----------------------------
\newcommand{\authorbio}[3]{%
    \noindent
    \begin{minipage}[t]{0.15\textwidth}
        \vspace{0pt}
        \centering
        \IfFileExists{#2}{%
            \includegraphics[width=0.95\linewidth,height=3.4cm,keepaspectratio]{#2}%
        }{%
            \relax%
        }
    \end{minipage}%
    \hfill
    \begin{minipage}[t]{0.81\textwidth}
        \vspace{0pt}
        \textbf{#1}\\[0.25em]
        #3
    \end{minipage}
    \par\vspace{0.6cm}
}

\begin{document}

% =============================================================================
% MASTHEAD
% =============================================================================
\noindent
\begin{minipage}[c]{0.16\textwidth}
    \centering
    \includegraphics[width=0.9\linewidth]{FOTOS/UTEC-Logo.jpg}
\end{minipage}%
\hfill
\begin{minipage}[c]{0.80\textwidth}
    \colorbox{mastheadgray}{
    \begin{minipage}{\linewidth}
        \centering
        \vspace{0.2cm}
        {\small Final Project 2026-2}\par
        \vspace{0.15cm}
        {\LARGE\bfseries Mechatronic Systems Design}\par
        \vspace{0.2cm}
        {\normalsize University of Engineering and Technology - UTEC}\par
        \vspace{0.1cm}
        {\small \url{www.utec.edu.pe}}\par
        \vspace{0.15cm}
    \end{minipage}
    }
\end{minipage}

\vspace{0.2cm}
\noindent\rule{\textwidth}{0.4pt}
\vspace{0.2cm}

% =============================================================================
% TITLE AND AUTHORS
% =============================================================================
\textit{Hardware article}

\vspace{0.25cm}

{\LARGE\bfseries A Vision-Based Ball-and-Plate Balancing Platform Mounted on a Mobile Robot}

\vspace{0.4cm}

\textbf{Angel Giovani Garcia Motta, Thomas Mat\'ias Flores Baltodano, Alejandro Guerra Figueroa, Carlos Eduardo Hilario Mayta, Alex Claudio Valenzuela Quispe}

\textit{Department of Mechatronic Engineering, University of Engineering and Technology (UTEC), Lima, Per\'u}

\vspace{0.4cm}

% =============================================================================
% ARTICLE INFO + ABSTRACT
% =============================================================================
\noindent
\begin{minipage}[t]{0.32\textwidth}
\textbf{ARTICLE INFO}\par
\vspace{0.05cm}
\noindent\rule{\linewidth}{0.4pt}
\vspace{0.05cm}
\textit{Keywords:}\\
Ball-and-plate system\\
Vision-based control\\
Linear quadratic regulator\\
Kalman filter\\
Mobile robot\\
NI myRIO\\
Raspberry Pi camera module\\
Servo-actuated platform
\end{minipage}%
\hfill
\begin{minipage}[t]{0.64\textwidth}
\textbf{ABSTRACT}\par
\vspace{0.05cm}
\noindent\rule{\linewidth}{0.4pt}
\vspace{0.05cm}
This article describes the design, construction and validation of a
ball-and-plate balancing platform mounted on a differential-drive mobile
robot. The base is driven by two DC motors with encoders and is
controlled by an NI myRIO with a proportional position law and a PD
wheel-speed loop selected among seven tested tunings. A $20\times20$~cm PLA plate on a central
spherical joint is tilted by two MG996R servomotors through
crank--connecting rod linkages. A Raspberry Pi 4B with an OV5647 camera
detects the 3~cm ball by HSV color segmentation, estimates its position
and velocity with a Kalman filter that compensates the camera delay,
and computes the plate tilt with an LQR controller with integral
action. The myRIO
enables the balancing loop through a digital line and receives the ball
position through UART. The sensing architecture was selected with a weighted Pugh matrix, and
the system was developed and validated following the V-model. In three runs on flat floor along
a C-shaped path, the mean distance between the ball and the plate
center was 0.96~cm and the ball stayed within
2~cm of the center 93\,\% of the time. In four runs on a $10^\circ$ ramp
these values were 1.29~cm and 78\,\%. The ball did not fall in any run.
\vspace{0.05cm}
\noindent\rule{\linewidth}{0.4pt}
\vspace{0.05cm}
\end{minipage}

\vspace{0.6cm}

% =============================================================================
% SPECIFICATIONS TABLE
% =============================================================================
\newpage

\textbf{Specifications table}\par

\begin{tabularx}{\textwidth}{|L{0.25\textwidth}|Y|}
\hline
\textbf{Hardware name} &
Mobile Robot with a Vision-Based Ball-and-Plate Balancing Platform \\
\hline
\textbf{Subject area} &
\begin{minipage}[t]{\linewidth}
\begin{itemize}[leftmargin=1.2em,nosep]
    \item Engineering and material science
    \item Educational tools and open-source alternatives to existing infrastructure
\end{itemize}
\end{minipage} \\
\hline
\textbf{Hardware type} &
\begin{minipage}[t]{\linewidth}
\begin{itemize}[leftmargin=1.2em,nosep]
    \item Electrical engineering and computer science
    \item Mechanical engineering and materials science
\end{itemize}
\end{minipage} \\
\hline
\textbf{Closest commercial analog} &
No commercial product was identified that combines a mobile robotic base
with a vision-based ball-and-plate mechanism for ball stabilization
during motion. Desktop ball-and-plate trainers exist for control
education, but they operate on a fixed base. \\
\hline
\textbf{Open-source license} &
GNU General Public License Version 2 (GPL-2.0). \\
\hline
\textbf{Cost of hardware} &
Approximately USD 1,398 (detailed in Table~\ref{tab:bom}). \\
\hline
\textbf{Source file repository} &
\url{https://github.com/alexsv315-lang/Vision-Based-Mobile-Ball-and-Plate-System.git} \\
\hline
\end{tabularx}

% =============================================================================
% MAIN BODY
% =============================================================================
\newpage

\section{Hardware in context}

Benchmark systems play a central role in control and mechatronics
because they allow algorithms and design decisions to be tested on a
physical plant whose behavior is well understood. Kaheman et al.\
illustrate this idea with the multi-arm pendulum on a cart, a benchmark
whose difficulty lies not only in its dynamics but also in the physical
constraints of the hardware, the control authority of the actuator, and
the placement of the sensors \cite{Kaheman2023}. The ball-and-plate
system, the two-dimensional extension of the ball-and-beam problem
\cite{Hauser1992}, plays a similar role: a ball must be kept near a
reference position by tilting a plate, which combines unstable,
nonlinear dynamics with the need for real-time sensing and actuation.
Awtar et al.\ presented one of the best-known implementations, a
ball-on-plate system designed and built by students in a mechatronics
course, documenting hardware design, sensor and actuator selection,
modeling, parameter identification, controller design and experimental
testing \cite{Awtar2002}. This makes the ball-and-plate a natural choice
for an educational mechatronic project.

Different sensing strategies have been used to measure the position of
the ball. Computer vision is attractive because it measures the ball
without physical contact with the plate. Park and Lee studied visual
servoing for the motion control of a ball on a plate \cite{Park2003},
and Cheng and Tsai developed a two-degree-of-freedom ball-and-plate
system in which a camera provided the ball position and a linear
quadratic regulator (LQR) commanded the plate orientation through two
servomotors \cite{ChengTsai2016}. Vision also has costs. Kaheman et al.\
point out that, compared with encoder-based measurements, vision-based
tracking requires additional computational resources, which is relevant
when high sampling rates are needed \cite{Kaheman2023}. In this journal,
Myers et al.\ showed that open-source video tracking can provide
high-quality data at low cost in a double-pendulum platform
\cite{Myers2020}. These works motivate the main architectural decision
of this project: to dedicate one processing unit to image acquisition,
ball-position estimation and plate control, so that the vision workload
does not compete with the locomotion control of the robot.

Most ball-and-plate implementations, however, are designed for a fixed
base. When the plate is carried by a moving vehicle, the motion of the
base and the slope of the terrain become additional disturbances that
the balancing system must reject. A related problem in robotics is the
waiter's problem, in which a robot must transport objects on a tray
without grasping them. Selvaggio et al.\ proposed a model predictive
non-sliding manipulation control for non-prehensile object
transportation \cite{Selvaggio2023}, Heins and Schoellig addressed the
problem on a mobile manipulator with a whole-body model predictive
controller that also avoids obstacles \cite{Heins2023}, and Brei et
al.\ studied real-time, safe motion planning and control for the
manipulation of unsecured objects \cite{Brei2024}. In these works the
transported object rests on a tray, the main concern is to prevent it
from sliding or tipping, and the platforms are mobile manipulators. In a
ball-and-plate system the ball rolls freely and the plate must be
actively tilted to keep it in place. The problem studied here combines
both ideas: a rolling object that must be kept near a reference position
while its supporting surface is carried by a moving base over flat and
inclined terrain.

Unlike conventional fixed-base ball-and-plate systems, the proposed
platform integrates active ball stabilization with differential-drive
locomotion. It therefore provides an experimental setup, built from
commercial educational parts, for evaluating ball-balancing control
under the disturbances caused by the motion of the vehicle and by
inclined terrain.

\section{Hardware description}

The hardware integrates two mechatronic subsystems on a single chassis
(Fig.~\ref{fig:proto_cad}): a differential-drive mobile base and a
two-axis ball-and-plate mechanism. A camera mounted on a mast above the
plate measures the position of the ball, and an inertial measurement
unit (IMU) on the plate is used to calibrate it. Fig.~\ref{fig:views}
shows the orthographic views of the CAD model and Fig.~\ref{fig:photos}
the built prototype on the bench and on the test track.

\begin{figure}[H]
    \centering
    \includegraphics[width=0.78\textwidth]{FOTOS/Fig01_prototype_CAD_annotated.png}
    \caption{CAD model of the complete prototype. Blue labels: mobile
    base subsystem; orange labels: balancing platform; purple labels:
    camera positioning system.}
    \label{fig:proto_cad}
\end{figure}

\begin{figure}[H]
    \centering
    \includegraphics[width=\textwidth]{FOTOS/Fig_views_CAD.png}
    \caption{Orthographic views of the CAD model: (a) front, (b) rear,
    (c) side and (d) top. The camera is mounted at the top of the mast,
    about 30~cm above the plate, looking down at its center.}
    \label{fig:views}
\end{figure}

\begin{figure}[H]
    \centering
    \includegraphics[width=\textwidth]{FOTOS/Fig05_real_prototype_photos.png}
    \caption{Built prototype. (a) Bench test of the balancing platform,
    (b) and (c) integrated robot running on the MDF test track with the
    ball balanced on the plate.}
    \label{fig:photos}
\end{figure}

The main design decisions and their justification are summarized below.
Each of them is developed in the following subsections.

\begin{itemize}
    \item \textbf{Distributed architecture with a self-contained balancing
    loop.} Locomotion runs on the NI myRIO and the complete balancing
    loop (camera, estimation, control and servo PWM) runs on the
    Raspberry Pi 4B. In the first version the servo PWM was generated by
    the myRIO, so every servo command had to cross the inter-processor
    link. Moving the servo outputs to a PCA9685 driver controlled by the
    Raspberry Pi removes that link from the loop: the UART latency only
    affects the position display, not the stability of the balancing
    loop.
    \item \textbf{Contact-free sensing with a fixed plate region.} A
    camera measures the ball without touching the plate
    \cite{Park2003,ChengTsai2016}. The four corners of the plate are
    fixed in the image instead of being tracked every frame
    (Section~\ref{sec:vision}).
    \item \textbf{Hue-based segmentation of a saturated ball on a white
    plate}, robust to the shading of the sphere and to changes of
    illumination (Section~\ref{sec:vision}).
    \item \textbf{Separate regulated supplies} for the Raspberry Pi and
    the servos, so that servo current peaks do not cause under-voltage
    of the computer (Section~\ref{sec:power}).
    \item \textbf{Two-axis platform driven by crank--connecting rod
    linkages}, a simple mechanism that converts the rotation of each servo
    into a tilt of the plate with a reduction ratio set by the geometry,
    which improves the angular resolution of the plate
    (Section~\ref{sec:platform}).
\end{itemize}

\subsection{Design methodology}
\label{sec:method}
The project was developed with the V-model for mechatronic systems
\cite{VDI2206}, shown in Fig.~\ref{fig:vmodel}. The requirements were
decomposed into a system design (two subsystems with separate
processors), a subsystem design (mechanism, power supply, vision,
control) and the implementation. Each design level was then verified by
a test at the same level: unit tests of each component, integration
tests of the platform with the base stopped, and system tests with the
robot moving on flat floor and on the $10^\circ$ ramp
(Section~\ref{sec:validation}). When a test did not meet its objective,
the corresponding design step was repeated; for example, the controller
was retuned from $\omega_n=5$ to $3$~rad/s after the integration tests.

\begin{figure}[H]
    \centering
    \includegraphics[width=\textwidth]{FOTOS/Fig00_V_model.png}
    \caption{V-model applied to the project. Each design step on the
    left is verified by the test at the same level on the right.}
    \label{fig:vmodel}
\end{figure}

The selection of the ball-position sensor, the component with the
largest effect on the architecture, was justified with a Pugh matrix
\cite{Pugh1991} (Table~\ref{tab:pugh}). A resistive touch panel, a
common sensor in ball-and-plate systems, was taken as the reference
(datum). Each alternative was scored as better ($+1$), equal ($0$) or
worse ($-1$) than the datum for each criterion, and the scores were
multiplied by the weight of the criterion (1 to 5), defined by the team
according to the requirements of a plate carried by a moving robot.

\begin{table}[H]
\centering
\caption{Pugh matrix for the selection of the ball-position sensor
(datum: resistive touch panel; $+1$ better, $0$ equal, $-1$ worse than
the datum).}
\label{tab:pugh}
\small
\begin{tabularx}{\textwidth}{@{}L{6.6cm}cYYY@{}}
\toprule
\textbf{Criterion} & \textbf{Weight} & \textbf{Touch panel (datum)} & \textbf{Camera + Raspberry Pi} & \textbf{Two IR distance sensors} \\
\midrule
No contact; any plate material (3D-printed PLA) & 4 & 0 & $+1$ & $+1$ \\
Reliable reading with a light 3~cm ball on a moving, vibrating base & 4 & 0 & $+1$ & $0$ \\
Both axes measured with a single sensor & 3 & 0 & $0$ & $-1$ \\
Sampling rate & 3 & 0 & $-1$ & $-1$ \\
Additional data for the interface (image, position, logging) & 3 & 0 & $+1$ & $0$ \\
Computational load & 2 & 0 & $-1$ & $0$ \\
Sensitivity to lighting & 2 & 0 & $-1$ & $-1$ \\
\midrule
\textbf{Weighted score} & & \textbf{0} & $\mathbf{+4}$ & $\mathbf{-4}$ \\
\bottomrule
\end{tabularx}
\end{table}

The camera obtained the highest score, mainly because it measures the
ball without contact and is not affected by the vibration of the moving
base. Its disadvantages (lower sampling rate, computational load and
sensitivity to lighting) were addressed in the design: a dedicated
processor, a reduced region of interest, a Kalman filter with delay
compensation, and HSV segmentation (Section~\ref{sec:vision}).

\subsection{Mobile base}
The mobile base has four wheels: two rear wheels driven by DC motors and
two passive omnidirectional wheels at the front. The chassis is built
with TETRIX aluminum channels (Fig.~\ref{fig:proto_cad}). Each driven
wheel has a radius of 51~mm and is powered by a 12~V TETRIX
TorqueNADO DC motor with a quadrature encoder read by the NI myRIO. The
motors are driven by a Cytron MDD10A dual H-bridge, which receives a PWM
signal and a direction signal per channel from the myRIO
\cite{MDD10A,myRIO}. Table~\ref{tab:base} summarizes the main
characteristics of the base. In tests carried out without the
ball-and-plate platform, the motors reached approximately 96~rpm, close
to the no-load speed given by the manufacturer.

\begin{table}[H]
\centering
\caption{Characteristics of the mobile base. Manufacturer values come
from the datasheet of the DC motors; the speed was measured by the
authors.}
\label{tab:base}
\small
\begin{tabularx}{\textwidth}{@{}L{6cm} Y@{}}
\toprule
\textbf{Parameter} & \textbf{Value} \\
\midrule
Wheels & 4 (2 driven, 2 omnidirectional) \\
Radius of the driven wheels & 51~mm \\
Motor supply voltage & 12~V \\
No-load speed (manufacturer) & 100~rpm \\
No-load current (manufacturer) & 200~mA \\
Stall torque (manufacturer) & 34~kg$\cdot$cm (470~oz$\cdot$in) \\
Stall current (manufacturer) & 5.5~A \\
Measured speed without platform load & approximately 96~rpm \\
Driver & MDD10A H-bridge, PWM and direction from the NI myRIO \\
Controller & NI myRIO-1900, LabVIEW \\
\bottomrule
\end{tabularx}
\end{table}

\subsection{Ball-and-plate platform}
\label{sec:platform}
The plate is a $20\times20$~cm PLA plate supported at its center by a
3D-printed column that ends in a spherical joint (Fig.~\ref{fig:mech}).
The joint carries the weight of the plate, so the servos only provide
the tilting moments. Each axis is actuated by one MG996R servomotor
through a crank--connecting rod linkage: the servo horn acts as the
crank, and a rod with articulated ends connects it to the underside of
the plate, at a distance $d$ from the central joint. When the crank of
radius $r$ rotates an angle $\varphi$ from the leveled position, the
end of the rod rises approximately $r\sin\varphi$ and the plate tilts
\begin{equation}
\theta \approx \arcsin\!\left(\frac{r}{d}\sin\varphi\right)
\approx \frac{r}{d}\,\varphi \quad\text{(small angles)}.
\label{eq:linkage}
\end{equation}
The ratio $G=r/d$ (degrees of plate per degree of servo) works as a
reduction: it reduces the torque required from the servo and improves
the angular resolution of the plate, at the cost of a smaller tilt
range. This mechanism was chosen because it is simple, it is the usual
solution for servo-driven ball-and-plate systems, and it can be built
with the servo horns, two rods and a 3D-printed support.

\begin{figure}[H]
    \centering
    \includegraphics[width=\textwidth]{FOTOS/Fig02_mechanism_CAD_annotated.png}
    \caption{Two-axis tilt mechanism. The plate pivots on a central
    spherical joint. Each axis is driven by one MG996R servomotor through
    a crank (servo horn) and a connecting rod attached to the underside
    of the plate.}
    \label{fig:mech}
\end{figure}

Since $r$, $d$ and the angles of the rods are not exactly the same on
both axes, $G$ was not calculated from the geometry but identified
experimentally by the control program: with no ball on the plate, each
servo is moved $\pm 6^\circ$ about its leveled position while the IMU
measures the change of inclination of the plate. The identified values
were
\begin{equation}
G_X = 0.21, \qquad G_Y = 0.30 ,
\end{equation}
with the X axis measured as roll and the Y axis as pitch of the IMU.
These gains determine the resolution and the range of the platform.
One count of the PCA9685 (12 bits at 50~Hz) is
$20\,000/4096 = 4.9~\mu$s, equivalent to $0.44^\circ$ of servo
rotation with the 500--2500~$\mu$s range used for $180^\circ$; on the
plate this is $0.09^\circ$ (X) and $0.13^\circ$ (Y). With the servo
travel limited by software to $\pm50^\circ$, the small-angle estimate
of the maximum plate tilt is $G\cdot 50^\circ = 10.5^\circ$ (X) and
$15^\circ$ (Y). Eq.~\eqref{eq:linkage} shows that the real range is
smaller, because the gain of the linkage decreases as the crank
rotates (it is proportional to $\cos\varphi$): with $\varphi=50^\circ$
it gives about $9.3^\circ$ (X) and $13.3^\circ$ (Y). The software
limit of the commanded tilt ($\pm10^\circ$, Section~\ref{sec:lqr}) is
common to both axes; on X the $\pm50^\circ$ servo clamp acts first and
limits the real tilt to about $9.3^\circ$. In all the logged runs the
largest commanded tilt about X was $4.4^\circ$, so this axis never
reached its mechanical limit. Separate limits per axis ($9^\circ$ on X,
$10^\circ$ on Y) are recommended for steeper slopes. This limitation is
discussed in Section~\ref{sec:validation}. Table~\ref{tab:platform}
summarizes the platform.

\begin{table}[H]
\centering
\caption{Characteristics of the ball-and-plate platform. Servomotor
values are taken from the manufacturer's datasheet \cite{MG996R}.}
\label{tab:platform}
\small
\begin{tabularx}{\textwidth}{@{}L{5.6cm} Y@{}}
\toprule
\textbf{Parameter} & \textbf{Value} \\
\midrule
Degrees of freedom & 2 (rotation about X and Y) \\
Plate & $20 \times 20$~cm, PLA (measured active area $19.45 \times 19.00$~cm) \\
Ball & Orange, 3~cm diameter \\
Central support & 3D-printed PLA column with spherical joint \\
Actuators & 2 $\times$ MG996R servomotors, one per axis \\
Transmission & Crank (servo horn) + connecting rod per axis \\
Identified gain plate/servo & $G_X=0.21$, $G_Y=0.30$ \\
Plate resolution (1 PWM count) & $0.09^\circ$ (X), $0.13^\circ$ (Y) \\
Servo supply & 7.0~V from an XL4016 buck converter \\
Servo operating voltage (datasheet) & 4.8 to 7.2~V \\
Servo stall torque (datasheet) & 9.4~kgf$\cdot$cm at 4.8~V; 11~kgf$\cdot$cm at 6~V \\
Servo command & PCA9685, 50~Hz, 500--2500~$\mu$s, 12-bit \\
Servo range used & $\pm 50^\circ$ about the leveled position \\
Leveled position (servo angle) & Set in the program; $\approx101^\circ$ (X) and $\approx97^\circ$ (Y) in the final assembly \\
Commanded plate tilt limit & $\pm 10^\circ$ \\
Plate orientation sensor & MPU-6500 IMU on the plate, I2C to the Raspberry Pi 4B \\
\bottomrule
\end{tabularx}
\end{table}

\subsection{Electrical architecture and power supply}
\label{sec:power}
Fig.~\ref{fig:wiring} shows the connections. The system has two
independent power domains with a common ground.

\begin{figure}[H]
    \centering
    \includegraphics[width=\textwidth]{FOTOS/Fig03_connection_diagram.png}
    \caption{Connection diagram. Red: power; blue: control signals; green:
    communication buses; purple: camera interface. The Raspberry Pi 4B
    drives the servos through the PCA9685 and reads the IMU on the same
    I2C bus. The myRIO enables the balancing loop with a digital line and
    receives the ball position through UART.}
    \label{fig:wiring}
\end{figure}

\textbf{Mobile base domain.} The 12~V NiMH TETRIX battery powers the
MDD10A H-bridge directly and the myRIO through an XL4016 buck converter
set to 10~V, within the 6--16~V input range of the myRIO-1900
\cite{myRIO}.
The encoders are powered at 5~V from the myRIO.

\textbf{Balancing platform domain.} A custom-assembled 12~V, 8~Ah
lithium-ion battery pack, with a nominal output of 12~V and a maximum
current of 5~A, feeds two buck converters. The cell configuration and
the protection circuit of the pack were not documented by its
assembler, and no fuse was used in the prototype. For a safe
reproduction, a pack with a BMS and a fuse of about 6~A in series with
its positive terminal are recommended, which also covers the worst-case
demand of the servos and the Raspberry Pi:

\begin{itemize}
    \item An \textbf{XL4015} (rated 5~A \cite{XL4015}) set to 5.1~V powers
    the Raspberry Pi 4B, which requires a 5.1~V, 3~A supply
    \cite{RPi4}. A dedicated converter is used because the Raspberry
    Pi reduces its clock or resets when its input drops below the
    under-voltage threshold. Servo current peaks on a shared supply
    would cause such drops. The converter works below 60\,\% of its
    rated current, which limits its heating during long continuous
    operation. The control program also logs under-voltage and
    temperature flags reported by the Raspberry Pi every 10~s, which
    allows this decision to be verified during the tests.
    \item An \textbf{XL4016} (rated 8~A \cite{XL4016}) set to 7.0~V powers
    the servos through the V+ terminal of the PCA9685. The two MG996R
    servos can draw about 2.5~A each at stall \cite{MG996R}, so the
    worst-case demand of 5~A is below the rating of the converter. The
    voltage is set at 7.0~V, inside the 4.8--7.2~V operating range,
    because the torque and speed of the servo increase with voltage. It
    must be adjusted with a multimeter before the servos are connected,
    since the maximum rating is only 0.2~V above.
\end{itemize}

\textbf{Signals.} The camera is connected to the CSI-2 port of the
Raspberry Pi with a 15-pin to 22-pin flat flex cable. The PCA9685
(address 0x40) and the MPU-6500 (address 0x68) share the I2C bus of the
Raspberry Pi and use 3.3~V logic. The PCA9685 generates the servo pulses
in hardware \cite{PCA9685}, which avoids the timing jitter of software
PWM under a non-real-time operating system.

\textbf{Link between the processors.} Two lines and a common ground
connect the Raspberry Pi and the connector A (MXP) of the myRIO
(Table~\ref{tab:link}). Both devices use 3.3~V logic \cite{myRIO}, so no
level shifter is needed.
\begin{itemize}
    \item \textbf{UART} (Raspberry Pi TX to myRIO RX): it sends the
    estimated ball position so that it can be plotted on the LabVIEW
    front panel. The format is 115200~baud, 8 data bits, no parity, one
    stop bit (8N1), without flow control. Each message is an ASCII line
    \texttt{x;y;e} terminated by a line feed (0x0A), where \texttt{x} and
    \texttt{y} are the position in cm with two decimals and \texttt{e}
    is the state of the control (1 = ON, 0 = OFF); for example,
    \texttt{-1.25;0.40;1}. Messages are sent at 20~Hz, which is enough
    for a plot. A text format was chosen because LabVIEW reads it
    directly with \textit{VISA Read} and \textit{Scan From String}, using
    the line feed as termination character. On the Raspberry Pi, the
    PL011 hardware UART is assigned to GPIO14/15 (serial console and
    Bluetooth disabled), because the alternative mini-UART changes its
    baud rate with the clock frequency of the processor \cite{RPiUART}. A
    loss of this link does not affect the balancing loop, which runs
    entirely on the Raspberry Pi.
    \item \textbf{Enable line} (myRIO DIO1 to Raspberry Pi GPIO27): a high
    level enables the balancing control and a low level disables it. The
    Raspberry Pi acts only on changes of level, with a 50~ms debounce, so
    the state of the line at start-up does not activate the control by
    itself. The DIO lines of the myRIO have internal 40~k$\Omega$ pull-up
    resistors to 3.3~V \cite{myRIO}; since no pull-down resistor was
    installed on GPIO27, the line goes high when the LabVIEW program is
    stopped. The start-up order of Section~\ref{sec:operation} avoids an
    unintended activation, and a 10~k$\Omega$ pull-down resistor is
    recommended as an improvement.
\end{itemize}

\begin{table}[H]
\centering
\caption{Connections between the Raspberry Pi 4B and the NI myRIO-1900
(MXP connector A).}
\label{tab:link}
\small
\begin{tabular}{@{}llll@{}}
\toprule
\textbf{Signal} & \textbf{Raspberry Pi 4B} & \textbf{NI myRIO (MXP A)} & \textbf{Settings} \\
\midrule
UART & GPIO14 / TXD (pin 8) & UART.RX (pin 10) & 115200 baud, 8N1, ASCII, LF \\
Enable & GPIO27 (pin 13), input & DIO1 (pin 13), output & high = ON, low = OFF \\
Ground & GND (pin 6) & DGND (pin 8) & common reference \\
\bottomrule
\end{tabular}
\end{table}

\subsection{Vision subsystem}
\label{sec:vision}
The ball position is measured with a Rev.~1.3 camera module for
Raspberry Pi based on the 5~MP OV5647 CMOS sensor
\cite{OV5647,RPiCamera}. The sensor supports up to 90 frames per
second at $640\times480$ \cite{OV5647}; the program configures
$640\times480$ at 60 frames per second, and the control loop runs at
the frame rate; the mean rate measured in the logs of the system tests
was 60.9 frames per second. Images are acquired with the Picamera2 library
\cite{Picamera2} in a dedicated thread that always keeps only the
newest frame, so the control loop never processes an old frame. If no
frame arrives for 2~s, the camera is restarted automatically.
The ball is detected with a \textbf{color-based blob detection
algorithm}: HSV color thresholding, morphological filtering, contour
extraction with the border-following algorithm of Suzuki and Abe
\cite{Suzuki1985}, shape filtering, and computation of the centroid with
image moments \cite{Gonzalez2018}. Fig.~\ref{fig:vision} shows the
processing chain, implemented with OpenCV \cite{OpenCV}.

\begin{figure}[H]
    \centering
    \includegraphics[width=\textwidth]{FOTOS/Fig07_vision_pipeline.png}
    \caption{Ball-position measurement pipeline running on the Raspberry
    Pi 4B.}
    \label{fig:vision}
\end{figure}

\textbf{Fixed plate region instead of plate tracking.} The program can
detect the white plate in every frame (threshold of low saturation and
high value, contour, polygon approximation and temporal median of the
corners). This option was discarded for operation for two reasons.
First, it requires processing the full $640\times480$ frame every cycle,
whereas with fixed corners only the region of interest around the plate
is processed, at half resolution (a quarter of its pixels). Second, in
preliminary tests the center of the tracked plate was not a fixed point
in the image: it varied by about 0.3~cm or more from frame to frame
because of the noise of the corner estimates. Since the ball position
is measured relative to that center, this variation entered the
controller directly as measurement noise. Instead, the four corners are
detected once with the plate leveled (or marked by hand with the mouse)
and saved in a file. A constant homography $H$ then maps image pixels to
plate coordinates in centimeters \cite{Hartley2004}. This assumption is
valid because the camera and the pivot of the plate are fixed to the
same chassis, so the center of the plate does not move in the image when
the plate tilts. The error introduced by the tilt grows with the
distance to the center and is zero at the center. At $10^\circ$ of
tilt, the projection of the edge of the plate (9.7~cm from the center)
shortens by $9.7(1-\cos 10^\circ)\approx0.15$~cm, and the edge rises
$9.7\sin 10^\circ \approx 1.7$~cm towards the camera. With the camera
$D\approx30$~cm above the plate, this produces a scale error of
$1.7/30\approx5.6\,\%$, about 0.55~cm at the edge. At 2~cm from the
center, where the ball is regulated, the same effects give less than
0.05~cm. The program includes an optional parallax correction that uses
the IMU angle and $D$; it was disabled in the reported tests because the
error near the reference is negligible.

\textbf{Why the ball appears blue to the camera.} With the
\texttt{RGB888} format, Picamera2 stores each pixel in memory in the
order [B,~G,~R] \cite{Picamera2}, while the program converts the image
to HSV with an RGB-to-HSV conversion. As a result, the red and blue
channels are swapped in everything the program sees. The real ball is
orange (OpenCV hue $\approx 5$, on a 0--179 scale), but in the swapped
image it has hue $\approx 115$, which is blue (Fig.~\ref{fig:color}b).
This is not a defect for the measurement: the threshold is defined and
tuned in the same swapped color space, so the segmentation is
consistent. The detection range is therefore H~=~100--130 in the
program, which corresponds to the orange of the real ball.

\textbf{Why HSV and these thresholds.} The ball is a sphere, so its
brightness is not uniform: the side facing the light is bright, the
opposite side is shaded, and the overall level changes with the
lighting of the room. In RGB space a fixed threshold mixes color and
brightness, so the darker parts of the ball are lost when the scene
darkens. HSV separates the chromaticity (hue H and saturation S) from
the brightness (value V) \cite{Gonzalez2018}. The program thresholds
mainly on hue (a narrow band of 30 units), requires a high saturation
($S\geq120$) to reject the white plate and specular highlights, which
have low saturation, and accepts a low value ($V\geq70$) so that the
shaded half of the ball is still detected. This is why a strongly
colored ball on a white plate was chosen: on the plate, the only
strongly saturated object is the ball. The thresholds can be adjusted
with sliders while viewing the binary mask. As Fig.~\ref{fig:color}c
shows, the color threshold alone also accepts other objects with a
similar color (the MDF floor of the track and some cables). For this
reason the mask is intersected with the fixed plate region
(Fig.~\ref{fig:color}d). Fig.~\ref{fig:interface} shows the interface
of the program on the Raspberry Pi during development.

\begin{figure}[H]
    \centering
    \includegraphics[width=\textwidth]{FOTOS/Fig09_interface.png}
    \caption{Program interface on the Raspberry Pi during development.
    (a) Binary mask of the ball after the HSV threshold. (b) Main window:
    plate contour (green), detected ball (red circle) with its centroid
    (blue) and position in cm (\textit{PLATO DETECTADO}: plate detected).
    (c) Mask of the white plate used by the plate-detection option; the
    hole corresponds to the ball.}
    \label{fig:interface}
\end{figure}

\begin{figure}[H]
    \centering
    \includegraphics[width=0.92\textwidth]{FOTOS/Fig08_color_detection.png}
    \caption{Color segmentation of the ball, reproduced offline on a frame
    of the test video with the same thresholds and filters as the
    program. (a) True colors, (b) image as interpreted by the program
    (R and B channels swapped by the RGB888 memory order), (c) HSV
    threshold alone, (d) threshold restricted to the fixed plate region
    and shape filters; the detected contour has a circularity of 0.90.}
    \label{fig:color}
\end{figure}

\textbf{Shape filters.} After a light morphological opening and closing
($3\times3$) and a median filter, each candidate contour must satisfy:
an equivalent diameter within $\pm35\,\%$ of the ball diameter (learned
automatically from the first detections), a circularity
$4\pi A/P^2 \geq 0.55$, a ratio between its area and that of its
minimum enclosing circle $\geq 0.55$, a centroid inside the plate, and
a distance to the previous detection below 5~cm. These filters reject a
hand or fingers entering the plate and isolated pixels. If the ball is
lost for a few frames, the last position is kept for up to 4 frames;
after 0.25~s without the ball, the plate returns to level.

\subsection{Plate orientation sensing}
The MPU-6500 combines a three-axis accelerometer and a three-axis
gyroscope; its datasheet specifies accelerometer full-scale ranges from
$\pm 2$~g to $\pm 16$~g and I2C communication at up to 400~kHz
\cite{MPU6500}. The IMU is fixed to the plate and is read at 200~Hz.
The inclination is obtained from the direction of gravity measured by
the accelerometer \cite{Pedley2013}. This estimate is correct only when
the sensor is not accelerating: any linear acceleration of the robot
adds to gravity and appears as a false inclination \cite{Pedley2013}.
The program checks the validity of every reading (frozen values,
communication errors, angles above $25^\circ$) and restarts the IMU
automatically when it fails. Its zero is stored with the plate leveled.
In the final configuration the IMU is used only for calibration: to
store the leveled position of the plate and to identify the gain $G$
(Section~\ref{sec:platform}). It is not used in the control loop, for
the reasons given in Section~\ref{sec:imu}.

\subsection{Design advantages}
\begin{itemize}
    \item \textbf{Modularity:} the mobile base and the balancing platform
    have separate processors and power supplies, so they can be tested
    and calibrated independently before integration.
    \item \textbf{Reproducibility:} the structure uses standard TETRIX
    parts and three simple 3D-printed parts, and all electronic modules
    are commercial.
    \item \textbf{Recovery routines:} the software includes automatic
    recovery routines for camera, IMU and I2C communication failures, and
    restarts the control loop after an unexpected error
    (Table~\ref{tab:faults}). These routines improve operational
    resilience, although they were not evaluated with systematic
    fault-injection tests.
    \item \textbf{Built-in data logging:} every control cycle is recorded
    (measured and estimated position, velocity, reference, commanded and
    measured tilt, servo angles, PWM and frame rate), and the results can
    be exported as plots for analysis.
\end{itemize}

\section{Control strategy}
\label{sec:control}

Fig.~\ref{fig:control} shows the two control loops, each organized as
reference, controller, actuator and plant, and sensor. They run on
different processors and are coupled only through the enable line and
the position display. The ball-and-plate loop uses only the camera as
feedback; the IMU is not part of the loop (Section~\ref{sec:imu}).

\begin{figure}[H]
    \centering
    \includegraphics[width=\textwidth]{FOTOS/Fig06_control_architecture.png}
    \caption{Control architecture, grouped as reference (yellow),
    controller (blue), actuator and plant (orange) and sensor (purple).
    (a) Ball-and-plate loop on the Raspberry Pi 4B: LQR with integral
    action on the camera measurement, one loop per axis. (b) Mobile base
    loop on the NI myRIO: proportional position law, differential-drive
    inverse kinematics and a PD speed loop per wheel.}
    \label{fig:control}
\end{figure}

\subsection{Ball model}
For small plate inclinations, a solid ball rolling without slipping on
the plate obeys, on each axis \cite{Awtar2002,Hauser1992},
\begin{equation}
\ddot{x} = \frac{5}{7}\, g \sin\theta \approx b\,\theta ,
\qquad b = \frac{5}{7}\,g\,\frac{\pi}{180} = 12.2~\mathrm{cm\,s^{-2}\,deg^{-1}},
\label{eq:ball}
\end{equation}
where $x$ is the ball position along the axis (cm), $g$ the gravity
acceleration (981~cm/s$^2$), $\theta$ the plate inclination (degrees)
and $b$ the gain between inclination and acceleration. The value of $b$
identified from the logged data was $8.0$~cm\,s$^{-2}$\,deg$^{-1}$,
lower than the theoretical value. The difference is attributed to
rolling friction and to the uncertainty of the gain $G$ that converts
plate angle into servo angle; the identified value was used for the
controller and the estimator. Defining the state
$\mathbf{z}=[e,\;v]^T$, where $e=x-r$ is the position error with
respect to the reference $r$ and $v=\dot x$ the velocity,
Eq.~\eqref{eq:ball} becomes a double integrator:
\begin{equation}
\dot{\mathbf{z}} = A\mathbf{z} + B\theta,\qquad
A=\begin{bmatrix}0&1\\0&0\end{bmatrix},\quad
B=\begin{bmatrix}0\\ b\end{bmatrix}.
\label{eq:ss}
\end{equation}

\subsection{LQR controller with integral action}
\label{sec:lqr}
The tilt is computed with the state-feedback law $\theta=-K\mathbf{z}$,
$K=[k_1,\;k_2]$, that minimizes the quadratic cost
\begin{equation}
J=\int_0^\infty\left(\mathbf{z}^TQ\mathbf{z}+R\,\theta^2\right)dt ,
\label{eq:cost}
\end{equation}
where $Q=\mathrm{diag}(q_1,q_2)$ weights the position and velocity
errors and $R$ the control effort \cite{Anderson1990}. For the double
integrator of Eq.~\eqref{eq:ss} with $R=1$, the algebraic Riccati
equation has the closed-form solution
\begin{equation}
k_1=\sqrt{q_1},\qquad
k_2=\sqrt{\frac{2\sqrt{q_1}}{b}+q_2} .
\label{eq:gains}
\end{equation}
The resulting closed loop is $\ddot e + b k_2\,\dot e + b k_1\,e = 0$,
with natural frequency $\omega_n=\sqrt{b k_1}$ and damping ratio
$\zeta=b k_2/(2\omega_n)$. Therefore, $q_1$ is chosen as
$(\omega_n^2/b)^2$ to place $\omega_n$ directly, and $q_2$ adds damping.

\textbf{Choice of the weights.} A preliminary tuning with
$\omega_n=5$~rad/s ($k_1=3.13^\circ$/cm, $k_2=0.94^\circ\cdot$s/cm) was
used in the bench tests of Section~\ref{sec:bench}. It produced fast
responses but large servo activity (Fig.~\ref{fig:servo}). For the final
tests the bandwidth was reduced to $\omega_n=3$~rad/s with $q_2=0.1$,
which gives
\begin{equation}
k_1=1.12~^\circ/\mathrm{cm},\qquad k_2=0.62~^\circ\cdot\mathrm{s/cm},
\qquad \zeta = 0.83,\qquad t_s\approx\frac{4}{\zeta\omega_n}=1.6~\mathrm{s},
\label{eq:final}
\end{equation}
where $t_s$ is the theoretical 2\,\% settling time. This choice
balances three limits of the hardware: the delay of the camera
(about 0.1~s, which reduces the phase margin at high bandwidth), the
limited tilt range of the linkage ($\approx9^\circ$ on X,
Section~\ref{sec:platform}), and the wear of the servos. With
$k_1=1.12^\circ$/cm, a 5~cm error requests $5.6^\circ$, within the
$\pm10^\circ$ limit, so the controller stays linear for most of the
plate. The damping $\zeta=0.83$ avoids large overshoots. The
implemented law adds four practical features:
\begin{enumerate}
    \item \textbf{Integral action} $-K_i\int e\,dt$ with
    $K_i=1.5~^\circ/(\mathrm{cm\cdot s})$, to remove steady-state errors
    due to leveling errors of the plate or to a slope. It is integrated
    only when $|e|<4$~cm and the output is not saturated, and it is
    frozen (not reset) otherwise. This conditional integration prevents
    windup \cite{Astrom2021}. Its contribution is limited to $3^\circ$.
    \item \textbf{Dead zones} of 0.4~cm on the error and 1~cm/s on the
    velocity, so the servos do not follow the noise of the measurement
    near the reference.
    \item \textbf{Limited approach speed}: the law is written as a
    velocity reference $v_d=-\mathrm{sat}(k_1e/k_2,\;30~\mathrm{cm/s})$
    and $\theta=-k_2(v-v_d)$, which is equivalent to the LQR law for
    small errors and avoids violent movements for large ones.
    \item \textbf{Saturation} of the commanded tilt at $\pm10^\circ$.
\end{enumerate}
With these features, the tilt commanded on each axis is
\begin{equation}
\theta_d=\mathrm{sat}_{\pm10^\circ}\!\left[-k_2\,(v_p-v_d)-\mathrm{sat}_{\pm3^\circ}\!\left(K_i\!\int e_{dz}\,dt\right)\right],
\qquad
v_d=-\mathrm{sat}_{\pm30}\!\left(\frac{k_1}{k_2}\,e_{dz}\right),
\label{eq:law}
\end{equation}
where $e_{dz}$ is the position error after the 0.4~cm dead zone and
$v_p$ the predicted velocity of Section~\ref{sec:kalman}. For small
errors Eq.~\eqref{eq:law} reduces to $\theta_d=-k_1e-k_2v-K_i\int e\,dt$.
The commanded servo angle is
\begin{equation}
\delta = \delta_0 + s\,\frac{\theta}{G},
\label{eq:servo}
\end{equation}
where $\delta_0$ is the leveled servo angle, $s=\pm1$ the direction of
the axis and $G$ the identified plate/servo gain. Its rate of change is
limited to $150^\circ$/s near the reference and $300^\circ$/s for large
errors, and changes below $0.3^\circ$ are not sent, to reduce wear.

\subsection{State estimation and delay compensation}
\label{sec:kalman}
The velocity is not measured; it is estimated with an independent
Kalman filter per axis \cite{Kalman1960,BarShalom2001} based on
Eq.~\eqref{eq:ball}, with the state $\mathbf{x}=[x,\;v]^T$, the
commanded tilt as input and the position measured by the camera as
output. With the sampling period $\Delta t$ of each frame
($\approx1/60$~s), the model and the filter are
\begin{equation}
\begin{gathered}
A_d=\begin{bmatrix}1&\Delta t\\0&1\end{bmatrix},\quad
B_d=\begin{bmatrix}\Delta t^2/2\\ \Delta t\end{bmatrix}b,\quad
H=\begin{bmatrix}1&0\end{bmatrix},\quad
Q_k=\sigma_a^2\begin{bmatrix}\Delta t^4/4&\Delta t^3/2\\ \Delta t^3/2&\Delta t^2\end{bmatrix},\quad
R_k=\sigma_m^2,\\[1mm]
\text{prediction:}\quad
\hat{\mathbf{x}}_{k|k-1}=A_d\hat{\mathbf{x}}_{k-1}+B_d\theta_{k-1},\qquad
P_{k|k-1}=A_dP_{k-1}A_d^T+Q_k,\\[1mm]
\text{update:}\quad
K_k=\frac{P_{k|k-1}H^T}{HP_{k|k-1}H^T+R_k},\quad
\hat{\mathbf{x}}_{k}=\hat{\mathbf{x}}_{k|k-1}+K_k\left(z_k-H\hat{\mathbf{x}}_{k|k-1}\right),\quad
P_k=(I-K_kH)P_{k|k-1},
\end{gathered}
\label{eq:kalman}
\end{equation}
where $z_k$ is the position measured by the camera, $P$ the error
covariance, $Q_k$ the covariance of the process noise, modeled as a
white acceleration of $\sigma_a=30$~cm/s$^2$, and $R_k$ the covariance
of the measurement noise, with $\sigma_m=0.15$~cm. The filter is
initialized when the ball is first detected, with
$\hat{\mathbf{x}}_0=[z_0,\;0]^T$ and
$P_0=\mathrm{diag}(1~\mathrm{cm}^2,\;100~\mathrm{cm^2/s^2})$.
Measurements that differ by more than 6~cm from the
prediction are rejected (gating); after five consecutive rejections
the filter is reinitialized with the new measurement. When the innovation exceeds 1.2~cm,
which indicates an external push, the process noise is multiplied by 8
for 0.3~s so that the filter adapts quickly. The acquisition and
processing of the image introduce a delay of about $\tau=0.1$~s. To
compensate it, the controller uses a prediction $\tau$ ahead,
$x_p=\hat x+\hat v\tau+\tfrac12 b\theta\tau^2$ and $v_p=\hat
v+b\theta\tau$, following the idea of predictor compensation of dead
time \cite{Smith1957}.

\subsection{Why the IMU is not used in the control loop}
\label{sec:imu}
An inner loop that corrected the servo command with the inclination
measured by the IMU was implemented and tested, but it is not used in
the final configuration, for three reasons. First, the IMU measures the
inclination of the plate with respect to gravity, which on the ramp
includes the $10^\circ$ of the terrain; cancelling it would require a
correction of about $10^\circ/G_Y\approx33^\circ$ of servo rotation on
Y, the direction of motion, far beyond the $5^\circ$ trim range of that
loop, which saturated and switched off every time the robot climbed the
ramp. Second, the
accelerations of the robot when it starts, stops, turns or enters the
ramp add to gravity in the accelerometer and are read as false
inclinations \cite{Pedley2013}. Third, the camera is fixed to the same
chassis as the plate, so it measures the ball directly in the frame of
the plate. Three angles must be distinguished along the direction of
motion: the slope of the terrain $\alpha$ (chassis with respect to the
horizontal), the tilt of the plate with respect to the chassis $\theta$
(the variable commanded by the controller through the servos) and the
tilt of the plate with respect to gravity, $\theta+\alpha$, which is
what the IMU measures and what accelerates the ball. Since the camera
reads positions in the plate frame and the controller acts on $\theta$,
the slope only appears as a constant disturbance $b\alpha$
in Eq.~\eqref{eq:ball}, and it is
compensated by the proportional and integral terms of the LQR. The
residual offset that remains on the ramp is analyzed in
Section~\ref{sec:ramp}.

\subsection{Mobile base control}
The mobile base is controlled with a LabVIEW program on the NI myRIO.
The encoders give the angular velocities of the wheels, from which the
pose of the robot is obtained by odometry \cite{Siegwart2011}. The
longitudinal position error is $e_x = x_d - x$, and the reference linear
velocity is
\begin{equation}
v_{\mathrm{ref}} =
\begin{cases}
\mathrm{sat}(K_p e_x), & |e_x| > \varepsilon \\[1mm]
0, & |e_x| \leq \varepsilon
\end{cases}
\end{equation}
with $K_p = 0.40$~s$^{-1}$ and $\varepsilon = 0.01$~m. The velocity limit
was set to $\pm 0.10$~m/s in the ball-balancing runs; a higher limit of
0.15~m/s was used only in the wheel-speed characterization of
Section~\ref{sec:wheel}. The low speed limit reduces the accelerations
transmitted to the ball. The angular velocity reference is zero during
longitudinal displacement. The differential-drive inverse kinematics
gives the wheel speed references \cite{Siegwart2011}:
\begin{equation}
\omega_R = \frac{v_{\mathrm{ref}} + \frac{L}{2}\omega_{\mathrm{ref}}}{r},
\qquad
\omega_L = \frac{v_{\mathrm{ref}} - \frac{L}{2}\omega_{\mathrm{ref}}}{r},
\end{equation}
where $L$ is the distance between the driving wheels and $r$ the wheel
radius, $L=0.33$~m and $r=0.051$~m. Each wheel speed is regulated
by a PD controller \cite{Astrom2021}. The LabVIEW PID block uses the
standard form
\begin{equation}
u(t)=K_c\left[e_\omega(t)+\frac{1}{T_i}\int e_\omega\,dt+T_d\,\frac{de_\omega}{dt}\right],
\label{eq:pidlv}
\end{equation}
with $T_i$ and $T_d$ entered in minutes in LabVIEW,
$e_\omega=\omega_{\mathrm{ref}}-\omega$ and $u$ the PWM duty cycle in
percent. The selected tuning (test PD-03, Section~\ref{sec:wheel}) uses
no integral action. Mathematically this corresponds to
$T_i\to\infty$; in the LabVIEW PID block it is configured by entering
$T_i=0$, a software convention that disables the integral term. The
implemented law is therefore the PD controller
\begin{equation}
u(t)=K_c\left[e_\omega(t)+T_d\,\frac{de_\omega}{dt}\right],
\qquad K_c=0.80~\%\cdot\mathrm{s/rad},\quad
T_d=0.010~\mathrm{min}=0.6~\mathrm{s},
\label{eq:pd}
\end{equation}
for both wheels. Its output is the PWM duty cycle and direction sent to the MDD10A. The feedback on wheel speed compensates the difference
between the two motors and the increase of load when the robot climbs
the ramp. The LabVIEW front panel shows the wheel
velocities, the pose and trajectory of the robot, and the position of
the ball received through UART.

\section{Design files summary}

The mechanical design combines commercial TETRIX components with three
custom 3D-printed parts: the plate, its central support and the
connecting rods. TETRIX was used for the chassis, the camera mast
and the links of the mechanism because its modular geometry allowed the
structure to be modified during development without permanent joints.
Additive manufacturing was used for these parts, whose geometry could
not be built with standard TETRIX parts.

\begin{table}[H]
\centering
\caption{Design files produced by the authors. Files marked ``GitHub'' are in the source file repository; the videos are on Google Drive.}
\label{tab:files_own}
\small
\begin{tabularx}{\textwidth}{|L{0.33\textwidth}|L{0.16\textwidth}|L{0.12\textwidth}|Y|}
\hline
\textbf{Design filename} & \textbf{File type} & \textbf{License} & \textbf{Location of the file} \\
\hline
TODITO.f3z (complete assembly) & Fusion 360 archive & GPL-2.0 & \href{https://github.com/alexsv315-lang/Vision-Based-Mobile-Ball-and-Plate-System}{GitHub} \\ \hline
BALL N PLATE, PLATE, CAMARA, BASEDSM, BASE POSTERIOR & Fusion 360 (inside the archive) & GPL-2.0 & \href{https://github.com/alexsv315-lang/Vision-Based-Mobile-Ball-and-Plate-System}{GitHub} \\ \hline
Central support of the plate & STL / STEP & GPL-2.0 & \href{https://github.com/alexsv315-lang/Vision-Based-Mobile-Ball-and-Plate-System}{GitHub} \\ \hline
39100\_TXM-Axle Spacers 1/8", 39101\_TXM-Axle Spacers & STEP & GPL-2.0 & \href{https://github.com/alexsv315-lang/Vision-Based-Mobile-Ball-and-Plate-System}{GitHub} \\ \hline
Control10-3\_v7\_UART.py & Python 3 & GPL-2.0 & \href{https://github.com/alexsv315-lang/Vision-Based-Mobile-Ball-and-Plate-System}{GitHub} \\ \hline
enlace\_myrio.py, mpu.py, exportar\_resultados.py & Python 3 & GPL-2.0 & \href{https://github.com/alexsv315-lang/Vision-Based-Mobile-Ball-and-Plate-System}{GitHub} \\ \hline
Mobile base program & LabVIEW VI & GPL-2.0 & \href{https://github.com/alexsv315-lang/Vision-Based-Mobile-Ball-and-Plate-System}{GitHub} \\ \hline
Connection diagram (Fig.~\ref{fig:wiring}) & PNG & GPL-2.0 & \href{https://github.com/alexsv315-lang/Vision-Based-Mobile-Ball-and-Plate-System}{GitHub} \\ \hline
log\_control\_10.csv (flat floor), log\_control\_11.csv (ramp) & CSV, semicolon-separated & GPL-2.0 & \href{https://github.com/alexsv315-lang/Vision-Based-Mobile-Ball-and-Plate-System}{GitHub} \\ \hline
analisis\_resultados.py (metrics and plots of Section~\ref{sec:validation}) & Python 3 & GPL-2.0 & \href{https://github.com/alexsv315-lang/Vision-Based-Mobile-Ball-and-Plate-System}{GitHub} \\ \hline
PD\_01, PD\_03\_P1--P3, PD2\_04\_1--3, PD2\_05\_1--3, PD2\_10, PD\_AGRESIVO, PID\_integral\_suave(\_en\_rampa) (wheel-speed tests) & CSV, comma-separated, 14 columns & GPL-2.0 & \href{https://github.com/alexsv315-lang/Vision-Based-Mobile-Ball-and-Plate-System}{GitHub} \\ \hline
Valores\_PID\_Datos\_\allowbreak correspondientes.pdf (gains of each test) & PDF & GPL-2.0 & \href{https://github.com/alexsv315-lang/Vision-Based-Mobile-Ball-and-Plate-System}{GitHub} \\ \hline
pid\_metrics.py, pid\_figs.py (metrics and plots of Section~\ref{sec:wheel}) & Python 3 & GPL-2.0 & \href{https://github.com/alexsv315-lang/Vision-Based-Mobile-Ball-and-Plate-System}{GitHub} \\ \hline
Video 1: ramp climb & MP4 & GPL-2.0 & \href{https://drive.google.com/file/d/12MyZ74ZD0Io3IhAB6Xtw1eezuV8RvqDj/view?usp=sharing}{Google Drive} \\ \hline
Video 2: reverse motion & MP4 & GPL-2.0 & \href{https://drive.google.com/file/d/1ND7cfqTfeYYMZg7W4a4FlZ_xd6O5o8VT/view?usp=sharing}{Google Drive} \\ \hline
\end{tabularx}
\end{table}

\begin{itemize}
    \item \textbf{TODITO.f3z} contains the complete assembly of
    Fig.~\ref{fig:proto_cad}, including the TETRIX parts, the
    electronic modules and the custom parts (plate, central support and
    connecting rods).
    \item \textbf{Control10-3\_v7\_UART.py} is the balancing program
    (vision, Kalman filter, LQR, inner loop, PCA9685 driver, logging). It
    needs \texttt{enlace\_myrio.py} (UART and enable line, Table~\ref{tab:link}),
    \texttt{mpu.py} (IMU driver) and \texttt{exportar\_resultados.py}
    (plots of the logged data) in the same folder.
    \item \textbf{log\_control\_10.csv} and \textbf{log\_control\_11.csv}
    are the raw logs of the tests of Section~\ref{sec:validation}, recorded
    by the balancing program at about 60 rows per second (time, measured
    and estimated position, velocity, reference, commanded tilt, servo
    angles, PWM and frame rate).
    \item The videos show the integrated robot climbing the $10^\circ$ ramp
    and moving in reverse with the ball balanced on the plate.
\end{itemize}

The TETRIX models used in the assembly are listed in
Table~\ref{tab:design_files}; they are proprietary and can be downloaded
from the manufacturer.

\begin{table}[H]
\centering
\caption{TETRIX design files used in the prototype.}
\label{tab:design_files}
\small
\begin{tabularx}{\textwidth}{|L{0.33\textwidth}|L{0.08\textwidth}|L{0.28\textwidth}|Y|}
\hline
\textbf{Design filename} & \textbf{Type} & \textbf{License} & \textbf{Location} \\
\hline
39068\_TXM-288mmChannel & STEP & Proprietary (Pitsco / TETRIX) & \href{https://www.pitsco.com/products/tetrix-channels}{TETRIX} \\ \hline
39073\_TXM-FlatBuilding Plates & STEP & Proprietary (Pitsco / TETRIX) & \href{https://www.pitsco.com/products/tetrix-flat-building-plates}{TETRIX} \\ \hline
39065\_TXM-32mmChannel & STEP & Proprietary (Pitsco / TETRIX) & \href{https://www.pitsco.com/products/tetrix-channels}{TETRIX} \\ \hline
39089\_TXM-MotorMount & STEP & Proprietary (Pitsco / TETRIX) & \href{https://www.pitsco.com/products/tetrix-motor-mount}{TETRIX} \\ \hline
44260\_TXM TorqueNADO DC Motor & STEP & Proprietary (Pitsco / TETRIX) & \href{https://www.pitsco.com/products/tetrix-max-torquenado-motor-with-encoder}{TETRIX} \\ \hline
39055\_TXM-4 inch Wheel & STEP & Proprietary (Pitsco / TETRIX) & \href{https://www.pitsco.com/products/tetrix-wheels}{TETRIX} \\ \hline
39079\_TXM-Motor Shaft Hub with Set Screw & STEP & Proprietary (Pitsco / TETRIX) & \href{https://www.pitsco.com/products/tetrix-motor-hub}{TETRIX} \\ \hline
39091\_TXM-Bronze Bushing & STEP & Proprietary (Pitsco / TETRIX) & \href{https://www.pitsco.com/products/tetrix-max-11-mm-bronze-bushing}{TETRIX} \\ \hline
36466\_TXM-4 inch Omni Wheel & STEP & Proprietary (Pitsco / TETRIX) & \href{https://www.pitsco.com/products/tetrix-omni-wheel-packs}{TETRIX} \\ \hline
39067\_TXM-160mmChannel & STEP & Proprietary (Pitsco / TETRIX) & \href{https://www.pitsco.com/products/tetrix-channels}{TETRIX} \\ \hline
39274\_TXM-Flat64mmx26mm & STEP & Proprietary (Pitsco / TETRIX) & \href{https://www.pitsco.com/products/tetrix-flats}{TETRIX} \\ \hline
39281\_TXM-InsideCornerBrackets & STEP & Proprietary (Pitsco / TETRIX) & \href{https://www.pitsco.com/products/tetrix-inside-corner-brackets}{TETRIX} \\ \hline
39273\_TXM-Flat96mmx26mm & STEP & Proprietary (Pitsco / TETRIX) & \href{https://www.pitsco.com/products/tetrix-flats}{TETRIX} \\ \hline
39070\_TXM-288mm Flat Bar & STEP & Proprietary (Pitsco / TETRIX) & \href{https://www.pitsco.com/products/tetrix-flat-bars}{TETRIX} \\ \hline
39066\_TXM-96mmChannel & STEP & Proprietary (Pitsco / TETRIX) & \href{https://www.pitsco.com/products/tetrix-channels}{TETRIX} \\ \hline
39272\_TXM-Flat160mmx26mm & STEP & Proprietary (Pitsco / TETRIX) & \href{https://www.pitsco.com/products/tetrix-flats}{TETRIX} \\ \hline
\end{tabularx}
\end{table}

All project files are also available in this
\href{https://drive.google.com/drive/folders/1bn01djJBaUqpOuXz7pNVUJY4Hojj-HjR?usp=sharing}{shared folder}.

\section{Bill of materials summary}

All costs are expressed in USD, rounded to whole dollars. Prices in
Peruvian soles were converted using an exchange rate of S/~3.39 per USD.
The TETRIX MAX structural and motion components are grouped into a
single line item corresponding to the complete TETRIX MAX set.

{\footnotesize
\begin{longtable}{|L{2.0cm}|L{3.6cm}|L{0.8cm}|L{1.6cm}|L{1.6cm}|L{2.6cm}|L{2.4cm}|}
\caption{Bill of materials.}\label{tab:bom}\\
\hline
\textbf{Designator} & \textbf{Component} & \textbf{Qty.} & \textbf{Cost per unit} & \textbf{Total cost} & \textbf{Source} & \textbf{Material type} \\
\hline
\endfirsthead
\hline
\textbf{Designator} & \textbf{Component} & \textbf{Qty.} & \textbf{Cost per unit} & \textbf{Total cost} & \textbf{Source} & \textbf{Material type} \\
\hline
\endhead
TETRIX MAX & TETRIX MAX complete set (structure, DC motors with encoders, wheels, 12~V NiMH battery) & 1 & USD 679 & USD 679 & Pitsco / TETRIX & Metal / Electromechanical \\ \hline
myRIO & NI myRIO-1900 embedded controller & 1 & USD 430 & USD 430 & National Instruments & Electronic \\ \hline
RPi 4B & Raspberry Pi 4 Model B (4~GB) & 1 & USD 144 & USD 144 & Raspberry Pi / supplier & Electronic \\ \hline
BAT2 & Custom 12~V, 8~Ah lithium-ion battery pack & 1 & USD 34 & USD 34 & Local supplier & Lithium-ion \\ \hline
MDD10A & Dual-channel 10~A DC motor driver & 1 & USD 26 & USD 26 & Cytron & Electronic \\ \hline
MG996R & Servomotor & 2 & USD 7 & USD 14 & Local supplier & Electromechanical \\ \hline
XL4016-S & 8~A buck converter module, set to 7.0~V (servos) & 1 & USD 7 & USD 7 & Local supplier & Electronic \\ \hline
XL4016-M & 8~A buck converter module, set to 10~V (myRIO) & 1 & USD 7 & USD 7 & Local supplier & Electronic \\ \hline
PLATE & Plate $20\times20$~cm (3D-printed) & 1 & USD 14 & USD 14 & 3D printing & PLA \\ \hline
PCA9685 & 16-channel PWM servo driver & 1 & USD 12 & USD 12 & Local supplier & Electronic \\ \hline
IMU & MPU-6500 inertial sensor (module sold as ``MPU-9250/6500'') & 1 & USD 10 & USD 10 & Local supplier & Electronic \\ \hline
CAM & Camera module Rev.~1.3 (OV5647) & 1 & USD 7 & USD 7 & Local supplier & Electronic \\ \hline
XL4015 & 5~A buck converter module (Raspberry Pi) & 1 & USD 6 & USD 6 & Local supplier & Electronic \\ \hline
FFC & CSI flat flex cable, 15 to 22 pins & 1 & USD 4 & USD 4 & Local supplier & Polymer / copper \\ \hline
SUP & Central support of the plate (3D-printed) & 1 & USD 3 & USD 3 & 3D printing & PLA \\ \hline
BALL & Orange ball, 3~cm diameter & 1 & USD 1 & USD 1 & Local supplier & Polymer \\ \hline
\multicolumn{4}{|r|}{\textbf{Total project cost}} & \multicolumn{3}{l|}{\textbf{USD 1,398}} \\
\hline
\end{longtable}
}
\noindent{\small Prices as of November 2026; local prices were converted from soles (PEN) and rounded to whole dollars. The TETRIX MAX set includes the structural parts, DC motors with encoders, wheels and 12~V NiMH battery listed in Table~\ref{tab:design_files}. Wires, connectors, jumper cables and fasteners outside the TETRIX set are not included in the total.}


\section{Build instructions}

\textbf{Tools and materials.} FDM 3D printer and PLA filament (plate,
central support and rods), TETRIX hex keys and wrenches, screwdrivers,
digital multimeter, soldering iron and heat-shrink tubing for the power
wiring, jumper wires for the logic signals, cable ties, and a computer
with Fusion 360, LabVIEW (with the myRIO module) and SSH access to the
Raspberry Pi. The files for each step are listed in
Tables~\ref{tab:files_own} and~\ref{tab:design_files}.

\textbf{Safety.} Disconnect both batteries while assembling and wiring.
Set the output voltage of the buck converters before connecting any
load. Keep fingers away from the cranks and rods when the servos are
powered: the servos can move suddenly when the program starts or when
the ball is detected.

\subsection{Mobile base}
\begin{enumerate}
    \item Assemble the rectangular chassis with TETRIX channels, flats and
    inside corner brackets, joined with screws and nuts
    (Table~\ref{tab:design_files}).
    \item Mount the two DC motors with their motor mounts at the rear,
    attach the 4-inch wheels with motor shaft hubs, and install the two
    omni wheels at the front with axles, bronze bushings and spacers.
    \item Build the camera mast with two vertical channels and a diagonal
    brace, and fix the camera bracket and the Raspberry Pi at the top
    (Fig.~\ref{fig:proto_cad}). The camera must look down at the center
    of the plate with the whole plate inside the image.
    \item Fix the myRIO, the MDD10A and the batteries on the chassis,
    keeping the center of mass low and between the axles.
\end{enumerate}

\subsection{Ball-and-plate mechanism}
\begin{enumerate}
    \item Print the plate, the central support and the connecting rods in
    PLA with 50\,\% infill and the default settings of a standard FDM
    slicer (0.2~mm layer height, two walls). Print each part on its
    largest flat face: the plate flat, with its top surface smoothed
    (sanded) so that the ball rolls without bumps; the rods lying flat;
    and the central support standing on its base. Fix the base of the central support to the chassis, below the
    center of the camera image, and attach the plate to its spherical
    joint so that it rotates freely in all directions.
    \item Fix the two servos to the chassis, one under the X axis and one
    under the Y axis of the plate (Fig.~\ref{fig:mech}).
    \item Before mounting the servo horns, power the servos and command
    the leveled positions defined in the program (about $101^\circ$ for
    X and $97^\circ$ for Y in the final assembly) with keys J/L and I/K of
    the program or with a servo tester. Mount the horns
    (cranks) in that position.
    \item Connect each crank to the underside of the plate with its
    connecting rod. Adjust the length of the rods so that the plate is
    level with the servos in the leveled position, and check that the
    articulated ends move without binding or noticeable clearance.
    \item Move each servo through its range and verify that the plate
    tilts without collisions between the rods, the plate and the chassis.
    \item Fix the IMU flat on the underside of the plate, aligned with
    the X and Y axes.
\end{enumerate}

\subsection{Electronics and wiring}
Follow Fig.~\ref{fig:wiring}.
\begin{enumerate}
    \item Connect the 12~V battery of the platform to the XL4015 and the
    XL4016. With no load, adjust the XL4015 to 5.1~V and the XL4016 to
    7.0~V with a multimeter.
    \item Connect the XL4015 output to the 5~V and GND pins of the
    Raspberry Pi, and the XL4016 output to the V+ terminal of the
    PCA9685.
    \item Connect the PCA9685 and the IMU to the I2C pins of the
    Raspberry Pi (SDA, SCL, 3.3~V, GND). Connect servo Y to channel 1
    and servo X to channel 2 of the PCA9685.
    \item Connect the camera to the CSI port with the flat flex cable
    (contacts facing the correct side of each connector).
    \item Connect the motors to the MDD10A, the MDD10A PWM/DIR inputs and
    the encoders to the myRIO, and the 12~V NiMH battery to the MDD10A and
    to the second XL4016, adjusted beforehand to 10~V, which powers the
    myRIO.
    \item Connect the UART line, the enable line and the common ground
    between the Raspberry Pi and the MXP connector A of the myRIO
    following Table~\ref{tab:link}. A 10~k$\Omega$ resistor from GPIO27 to
    GND is recommended (Section~\ref{sec:power}).
    \item On the Raspberry Pi, run \texttt{sudo raspi-config} and, in
    \textit{Interface Options / Serial Port}, disable the login shell and
    enable the serial hardware; add \texttt{dtoverlay=disable-bt} to
    \texttt{/boot/firmware/config.txt} and reboot.
    \item In the LabVIEW program of the myRIO, configure the UART of
    connector A (VISA resource of connector A) at 115200~baud, 8 data bits,
    no parity, one stop bit, with the termination character enabled and
    set to 0x0A. In the loop, read a line with \textit{VISA Read} and
    convert it with \textit{Scan From String} using the format
    \texttt{\%.;\%f;\%f;\%d} (the code \texttt{\%.;} forces the point as
    decimal separator). Write the enable button to connector A DIO1 with a
    digital output.
\end{enumerate}

\subsection{Verification}
Before operation, check the items of Table~\ref{tab:checks}.

\begin{table}[H]
\centering
\caption{Verification checklist after assembly.}
\label{tab:checks}
\small
\begin{tabularx}{\textwidth}{@{}L{5.2cm}Y@{}}
\toprule
\textbf{Check} & \textbf{Expected result} \\
\midrule
Converter outputs (no load) & XL4015: 5.1~V; XL4016 (servos): 7.0~V; XL4016 (myRIO): 10~V \\
Mechanism & Plate level with the servos at their leveled positions; free tilt without collisions \\
Wheels and motors & Free rotation; both wheels turn forward with a positive command from LabVIEW \\
Camera & Whole plate visible in the image; ball detected (key M shows a single round blob) \\
IMU & Inclination close to $0^\circ$ with the plate leveled \\
UART link & Ball position appears on the LabVIEW front panel \\
Enable line & The program shows CONTROL: ON when the LabVIEW button is switched ON \\
\bottomrule
\end{tabularx}
\end{table}

\section{Operation instructions}
\label{sec:operation}

\subsection{Software installation}
On the Raspberry Pi (Raspberry Pi OS), enable the camera and the I2C
interface, and install Python~3 with \texttt{picamera2}, \texttt{opencv},
\texttt{numpy}, \texttt{smbus2}, \texttt{pyserial} and
\texttt{gpiozero}. The experiments used a Raspberry Pi 4 Model B (4~GB) with Raspberry Pi
OS and the versions of Python~3, OpenCV and Picamera2 available in its
package repositories in November 2026, and LabVIEW with the myRIO
module on the NI myRIO. Since exact version numbers were not recorded,
users should record them with \texttt{pip list} and keep them fixed when
reproducing the results. Copy the four Python files of
Table~\ref{tab:files_own} to the same folder. Load the LabVIEW program
into the myRIO.

\subsection{Start-up and calibration}
\begin{enumerate}
    \item Place the robot on a level surface, power the mobile base and
    then the platform.
    \item Start the LabVIEW program first, with the enable button in
    OFF. Because GPIO27 has no pull-down resistor, starting the Raspberry
    Pi program while the myRIO output is not driven could enable the
    control unintentionally.
    \item Run \texttt{python3 Control10-3\_v7\_UART.py}. Enter the
    reference position (Enter twice selects the center of the plate).
    \item The servos move slowly to the leveled position. The first time
    (or if the calibration files are missing), the program stores the
    zero of the IMU and identifies the gains $G$ automatically. The
    plate must be empty during this step (about 15~s).
    \item The program fixes the plate region automatically when its
    corners are stable, or loads the corners saved in
    \texttt{plato\_fijo.json}. If the plate is not well framed, press B
    and click the four corners. Verify the detection of the ball with
    key M (binary mask) and adjust the HSV sliders if the lighting
    changes.
\end{enumerate}

\subsection{Operation}
\begin{enumerate}
    \item Place the ball on the plate and enable the control with the
    myRIO digital line (ON), or with key C. The display shows the state,
    the frame rate, the measured and estimated position, the servo
    commands and the IMU state.
    \item The reference can be changed with a click on the image or with
    keys W/A/S/D (1~cm steps). Keys 1--8 change the bandwidth, damping,
    $K_i$ and dead zone during operation; P prints all parameters.
    \item Start the motion of the mobile base from the LabVIEW front panel.
    The ball position received through UART is plotted on the panel.
\end{enumerate}

\subsection{Shutdown}
Disable the control from LabVIEW (button OFF) or with key C: the plate
returns to level. Press Q to exit the Raspberry Pi program, and only
then stop the LabVIEW program; the program levels the plate and exports the
logged data and plots to \texttt{resultados/<date\_time>/}. Stop the
mobile base, then disconnect the platform battery and the base battery.

\textbf{Reading the results.} The window of the program shows the
state of the control, the frame rate, the measured and estimated
position (cm), the velocity (cm/s), the servo angles and PWM ($\mu$s)
and the LQR gains. The LabVIEW front panel shows the pose and velocity
of the robot and the ball position received through UART. After each
session, the log \texttt{log\_control.csv} and the plots exported to
\texttt{resultados/} can be analyzed with the metrics of
Section~\ref{sec:validation}.

\textbf{Behavior under faults.} Table~\ref{tab:faults} summarizes how
the software reacts to the most common faults. These reactions are
implemented in the code but were not tested systematically.

\begin{table}[H]
\centering
\caption{Reaction of the system to faults (as implemented in the code).}
\label{tab:faults}
\small
\begin{tabularx}{\textwidth}{@{}L{4.4cm}Y@{}}
\toprule
\textbf{Fault} & \textbf{Reaction} \\
\midrule
Ball not detected in a frame & The last position is kept for up to 4 frames; after 0.25~s without the ball the plate returns to level \\
Outlier measurement (jump $>6$~cm) & Rejected by the Kalman gating; after 5 consecutive rejections the filter restarts at the new position \\
UART link lost & A warning is printed and the balancing loop continues, since it runs entirely on the Raspberry Pi; only the plot in LabVIEW is lost \\
Enable line set to OFF & The control is disabled and the plate returns to level \\
Camera, IMU or I2C error & The program retries the communication automatically; I2C errors are counted and reported \\
Restart of the myRIO & GPIO27 floats without a pull-down; follow the start-up order of this section \\
\bottomrule
\end{tabularx}
\end{table}

\textbf{Operating limits.} Table~\ref{tab:limits} lists the conditions
under which the prototype was validated.

\begin{table}[H]
\centering
\caption{Operating limits of the prototype.}
\label{tab:limits}
\small
\begin{tabularx}{\textwidth}{@{}L{5.2cm}Y@{}}
\toprule
\textbf{Parameter} & \textbf{Limit} \\
\midrule
Terrain slope & Up to $10^\circ$ (tested); the plate tilt range ($\approx9^\circ$ on X) limits larger slopes \\
Base speed & $\leq0.10$~m/s (software limit) \\
Ball & Orange, 3~cm diameter; another color requires retuning the HSV thresholds \\
Lighting & Indoor, uniform light; recalibrate HSV if the lighting changes \\
Reference & At least 2~cm from the edges of the plate \\
Supplies & Do not exceed 7.2~V on the servos; recharge the batteries before long tests \\
\bottomrule
\end{tabularx}
\end{table}

\textbf{Hazards.} Pinch points at the cranks and rods; possible sudden
servo motion at start-up; loss of the ball from the plate if the base
accelerates abruptly; heating of the converters during long operation.
Do not leave the batteries connected unattended.

\section{Validation and characterization}
\label{sec:validation}

\subsection{Test plan, metrics and acceptance criteria}
Following the V-model (Fig.~\ref{fig:vmodel}), the prototype was
validated at three levels: unit tests of the wheel-speed controller of
the base, integration tests of the platform with the base stopped, and system tests with the robot moving on flat floor and
on the $10^\circ$ ramp (Table~\ref{tab:plan}). In the platform tests the
reference was the center of the plate, $(0,0)$, and the Raspberry Pi
logged every control cycle (about 60 rows per second).

\begin{table}[H]
\centering
\caption{Test plan.}
\label{tab:plan}
\small
\begin{tabularx}{\textwidth}{@{}L{3.2cm}L{1.0cm}YL{3.0cm}@{}}
\toprule
\textbf{Test} & \textbf{Runs} & \textbf{Condition} & \textbf{Controller under test} \\
\midrule
Unit (wheel-speed loop) & 14 & Base with the platform, 3~m out-and-back route on flat floor; one run on the $10^\circ$ ramp & Seven PD/PID tunings of the wheel speed \\
Integration (bench) & 5 & Base stopped; ball released 5.7--8.0~cm from the center & Preliminary, $\omega_n=5$~rad/s \\
System, flat floor & 3 & Robot moving along a C-shaped path (straight sections and turns) and returning to its start & Final, $\omega_n=3$~rad/s \\
System, ramp & 4 & Robot climbing and descending the $10^\circ$ ramp, then reversing to its start & Final, $\omega_n=3$~rad/s \\
\bottomrule
\end{tabularx}
\end{table}

For each run, the error on each axis was computed as $e_x=\hat x-r_x$
and $e_y=\hat y-r_y$, where $\hat x,\hat y$ are the estimated positions
and $r_x, r_y$ the reference, and the distance to the reference as
$d=\sqrt{e_x^2+e_y^2}$. The metrics are:
\begin{itemize}
    \item \textbf{Mean error per axis}: mean of $|e_x|$ and of $|e_y|$.
    \item \textbf{Mean distance}: mean of $d$; it summarizes in one number
    how close to the reference the ball stayed.
    \item \textbf{Maximum distance}: largest value of $d$ during the run.
    \item \textbf{Time within 2~cm}: percentage of samples with
    $d<2$~cm. It is used as the \textit{effectiveness} of the regulation.
\end{itemize}
Distances are also expressed as a percentage of the plate side (20~cm).
Two acceptance criteria were defined by the team: (1) the ball must not
fall from the plate in any run, and (2) the mean distance must be below
2~cm, that is, 10\,\% of the plate side.

\textbf{Data processing.} The logs of the system tests
(\texttt{log\_control\_10.csv} and \texttt{log\_control\_11.csv},
Table~\ref{tab:files_own}) were processed as follows: (1) sessions were
separated where the time column restarts, because the program was
restarted between some runs; (2) each run was defined as a continuous
period with the ball on the plate, delimited by gaps longer than 1~s
without detection, which correspond to the ball being placed or removed
by hand between runs; isolated detections of a few rows outside these
periods (for example, while the ball was being held by hand) were not
considered runs; (3) 2~s were removed at the start and end of each run;
(4) in flat run 2, the last 4 minutes, with the robot already stopped
and the ball still inside the dead zone, were removed so as not to
improve the result artificially. To check that this processing does
not bias the metrics, the availability of the vision system was
computed inside each run: the ball was detected in 100\,\% of the
frames of all seven runs, and the longest interval between two
consecutive detections was 0.03~s (two frames). Therefore, no samples
were removed inside the runs, and the metrics include every instant of
the regulation. The duration of each run is given in
Tables~\ref{tab:flat} and~\ref{tab:ramp}. All the tests were carried
out with a single prototype and a single 3~cm ball.

\subsection{Unit tests: wheel-speed controller of the mobile base}
\label{sec:wheel}
\textbf{Procedure.} The base, carrying the platform, followed the same
route in every test: a 3~m straight segment at $v_{\mathrm{ref}}\le
0.15$~m/s (above the 0.10~m/s of the balancing runs, to test the loop
with a larger margin) ($\omega_{\mathrm{ref}}=2.94$~rad/s per wheel), a stop, a
$180^\circ$ turn in place at 0.3~rad/s and the return to the origin
(Fig.~\ref{fig:basetraj}). The myRIO logged 14 variables at about 92~Hz: elapsed time; reference
speed, measured speed and PWM of the right wheel
(\texttt{wd\_ref}, \texttt{D\_Rad/s}, \texttt{PWM\_D}) and of the left
wheel (\texttt{wi\_ref}, \texttt{I\_Rad/s}, \texttt{PWM\_I}); the
odometric pose (\texttt{x}, \texttt{y}, \texttt{theta}); the velocity
references (\texttt{v\_ref}, \texttt{omega\_ref}); and the distance and
heading errors (\texttt{rho}, \texttt{e\_theta}) (\texttt{PD\_*.csv}, \texttt{PD2\_*.csv}
and \texttt{PID\_*.csv}, Table~\ref{tab:files_own}). Seven tunings of
Eq.~\eqref{eq:pidlv} were tested, first with equal gains on both wheels
($K_c$ from 0.60 to 1.10), then with different gains per wheel to
compensate the asymmetry of the motors, and finally with a soft integral
action ($T_i=1.5$~min $=90$~s). The tuning in use (PD-03) and the asymmetric tunings PD2-04 and
PD2-05 were repeated three times.

\textbf{Metrics.} During cruise ($v_{\mathrm{ref}}=0.15$~m/s, after
the first 3~s) the following were computed for each wheel: the mean
absolute speed error as a percentage of the reference, and the standard
deviation of the PWM (ripple, an indicator of control effort and motor
wear). The maximum PWM and the percentage of saturated samples
($u\ge 99\,\%$) were computed over the whole run. The heading
deviation is the odometric heading $|\theta|$ at the end of the 3~m
straight segment; it measures how well the two wheels are matched,
which is what keeps the platform on a straight path.

\begin{table}[H]
\centering
\caption{Comparison of the wheel-speed tunings (mean of the runs;
$\pm$: standard deviation of 3 runs). $K_c$ is given for the right /
left wheel; $T_d=0.010$~min unless noted.}
\label{tab:wheel}
\small
\begin{tabularx}{\textwidth}{@{}L{2.0cm}L{3.1cm}cYYYc@{}}
\toprule
\textbf{Test} & \textbf{Gains} & \textbf{Runs} & \textbf{Speed error [\%]} & \textbf{Heading dev. [$^\circ$]} & \textbf{PWM ripple [\%]} & \textbf{Max PWM [\%]} \\
\midrule
PD-01 & $K_c$ 0.60 / 0.60 & 1 & 5.2 & 0.11 & 1.16 & 32 \\
\textbf{PD-03 (selected)} & $K_c$ 0.80 / 0.80 & 3 & 5.2 & \textbf{0.06 $\pm$ 0.05} & \textbf{1.10} & 36 \\
PD2-04 & $K_c$ 0.70 / 0.80 & 3 & 5.2 & 0.22 $\pm$ 0.12 & 1.15 & 34 \\
PD2-05 & $K_c$ 0.80 / 0.70 & 3 & 5.2 & 0.11 $\pm$ 0.08 & 1.10 & 36 \\
PD2-10 & $K_c$ 2.0 / 0.90, $T_d$ 0.008 / 0.012 & 1 & 5.2 & 0.27 & 1.17 & 38 \\
PD aggressive & $K_c$ 1.10 / 0.95, $T_d$ 0.015 & 1 & 5.1 & 0.12 & 1.16 & 34 \\
PID soft & $K_c$ 0.80, $T_i$ 1.5~min & 1 & 5.0 & 0.11 & 1.80 & 100 \\
PID soft, ramp & $K_c$ 0.80, $T_i$ 1.5~min & 1 & 5.3 & -- & 3.39 & 100 \\
\bottomrule
\end{tabularx}
\end{table}

\begin{figure}[H]
    \centering
    \includegraphics[width=\textwidth]{FOTOS/Fig19_base_gains.png}
    \caption{Comparison of the wheel-speed tunings on flat floor.
    (a) Mean speed error during cruise. (b) Heading deviation at the end
    of the 3~m straight segment. (c) PWM ripple. The selected tuning,
    PD-03, is shown in green.}
    \label{fig:basegains}
\end{figure}

\textbf{Results.} The speed error was between 5.0 and 5.3\,\% for all
the tunings (Table~\ref{tab:wheel}). This value is set by the encoder,
not by the controller: at the 92~Hz sampling rate one encoder count
corresponds to 0.218~rad/s, that is, 7.4\,\% of the 2.94~rad/s
reference, so the measured speed jumps between neighboring levels
(Fig.~\ref{fig:basetraj}b) and differences smaller than one count cannot
be resolved. For the same reason, the overshoot at start-up (11, 19 or
26\,\%) was always one, two or three counts. The tunings were therefore
compared by the heading deviation and by the control effort
(Fig.~\ref{fig:basegains}):
\begin{itemize}
    \item \textbf{PD-03} ($K_c=0.80$, $T_d=0.010$~min on both wheels) gave
    the smallest heading deviation ($0.06^\circ$ mean over three runs) and
    the lowest PWM ripple (1.10\,\%), without saturation. It was selected.
    \item \textbf{Different gains per wheel} (PD2-04, PD2-05, PD2-10) did
    not improve the straightness: the deviation increased to
    0.11--$0.27^\circ$ and was less repeatable. The speed loop already
    compensates the difference between the motors.
    \item \textbf{Higher gains} (PD aggressive) did not reduce the error,
    because it is limited by the encoder resolution.
    \item \textbf{Integral action} (PID soft) did not reduce the error
    either, increased the PWM ripple by 64\,\% and produced 100\,\% PWM peaks at start-up and when the return began,
    caused by the integral term accumulated while the robot was stopped.
\end{itemize}

\begin{figure}[H]
    \centering
    \includegraphics[width=\textwidth]{FOTOS/Fig18_base_trajectory.png}
    \caption{Base trajectory with the selected controller (PD-03, run 1).
    (a) Odometric position $x$ and heading $\theta$: 3~m forward,
    $180^\circ$ turn and return to the origin. (b) Reference and measured
    speed of both wheels; the steps of the measured speed are the
    encoder resolution. (c) PWM of both wheels (control effort).}
    \label{fig:basetraj}
\end{figure}

Fig.~\ref{fig:basetraj} shows the complete route with PD-03. The
odometric error at the end of every run was 0.01~m, which is the
stopping tolerance of the program. The real error, measured by hand
between the final and the initial position of the robot in five runs,
was 1.5, 1.0, 0.9, 1.0 and 1.4~cm, that is, $1.16 \pm 0.27$~cm (mean
$\pm$ s.d.), 0.4\,\% of the 6~m route. The difference with the
odometry is due to wheel slip and to the encoder resolution. The PWM stayed between 27 and 30\,\% during
cruise and the brief peak at the start of the return (34\,\%) is far
from saturation.

\textbf{Effect of the ramp.} Only the PID soft tuning was logged on the
ramp. Since it was the same route, it gives the extra effort needed to
climb (Fig.~\ref{fig:baseramp}): during the climb the mean PWM rose from
29.5 to 35.6\,\% on the right wheel and from 27.9 to 32.6\,\% on the
left (+17 to +21\,\%), and during the descent it fell by 7 to 9\,\%. The
speed error during the climb stayed at 4.7\,\% (5.5\,\% when climbing in
reverse), equal to that on flat floor, so the speed loop compensates the
load of the slope. The PWM reached during the climb (about 37\,\%) leaves a
large margin below saturation. The PD-03 tuning, used in all the
platform tests of the following sections including the ramp runs, was
not characterized separately on the ramp. Since it has the same
proportional and derivative gains, a similar compensation of the slope
is expected, but this is an inference from the gains, not a measured
result.

\begin{figure}[H]
    \centering
    \includegraphics[width=\textwidth]{FOTOS/Fig20_base_ramp.png}
    \caption{PWM of the wheels on flat floor and on the $10^\circ$ ramp
    with the same route and controller (moving average of 15 samples).
    Orange bands: climb; blue bands: descent.}
    \label{fig:baseramp}
\end{figure}

\subsection{Integration tests: platform with the base stopped}
\label{sec:bench}
Five regulation tests were carried out with the preliminary tuning. In
each one the ball was released at a different initial position, in
different quadrants. Two additional metrics were computed: the time to
first reach the reference, and the maximum distance after that moment.
Table~\ref{tab:results} summarizes the results, Fig.~\ref{fig:traj}
shows the trajectories, and Fig.~\ref{fig:time} the position as a
function of time.

\begin{table}[H]
\centering
\caption{Integration tests with the base stopped (preliminary tuning,
reference at the center of the plate).}
\label{tab:results}
\footnotesize
\setlength{\tabcolsep}{4pt}
\begin{tabular}{@{}lccccccc@{}}
\toprule
\textbf{Run} & \textbf{Initial position} & \textbf{Initial} & \textbf{Time to} & \textbf{Max.\ dist.} & \textbf{Max.\ dist.} & \textbf{Mean dist.} & \textbf{Mean dist.} \\
 & $(x_0,y_0)$ [cm] & \textbf{distance} [cm] & \textbf{reference} [s] & [cm] & [\% of plate] & [cm] & [\% of plate] \\
\midrule
1 & $(-5.63,\;0.63)$  & 5.67 & 0.60 & 3.115 & 15.57 & 0.799 & 3.99 \\
2 & $(5.75,\;5.53)$   & 7.98 & 2.32 & 0.528 & 2.64  & 0.226 & 1.13 \\
3 & $(4.60,\;-3.56)$  & 5.82 & 0.62 & 2.877 & 14.39 & 0.844 & 4.22 \\
4 & $(-5.42,\;-4.25)$ & 6.89 & 2.09 & 0.824 & 4.12  & 0.493 & 2.47 \\
5 & $(0.14,\;6.64)$   & 6.64 & 0.59 & 0.955 & 4.78  & 0.428 & 2.14 \\
\midrule
Mean $\pm$ s.d. & -- & -- & $1.24\pm0.88$ & $1.66\pm1.23$ & $8.3\pm6.2$ & $0.56\pm0.26$ & $2.8\pm1.3$ \\
\bottomrule
\end{tabular}
\end{table}

\begin{figure}[H]
    \centering
    \includegraphics[width=\textwidth]{FOTOS/Fig10_trajectories_5tests.png}
    \caption{Trajectory of the ball on the plate in the five integration
    tests. Green dot: start (\textit{Inicio}); red square: end
    (\textit{Fin}); black cross: reference (\textit{Referencia}); blue
    line: Kalman estimate (\textit{Estimada}); gray dots: camera
    measurement (\textit{Medida}).}
    \label{fig:traj}
\end{figure}

\begin{figure}[H]
    \centering
    \includegraphics[width=\textwidth]{FOTOS/Fig11_time_response_5tests.png}
    \caption{Estimated position $X(t)$ and $Y(t)$ of the ball (blue) and
    reference (dashed) in the five integration tests. Time in seconds
    (\textit{Tiempo}).}
    \label{fig:time}
\end{figure}

In all five tests the ball was brought to the reference, with a mean
distance of $0.56$~cm (2.8\,\% of the plate side). In runs 1, 3 and 5
the ball reached the reference in about 0.6~s but then overshot to about
3~cm; in runs 2 and 4, which started farthest from the center, the tilt
saturated and the ball oscillated before settling at 2.1--2.3~s. These
tests also revealed a limitation of the preliminary tuning: after the
transient, the servos kept oscillating with an amplitude of
20--30$^\circ$ while the ball stayed within $\pm1$~cm
(Fig.~\ref{fig:servo}). This limit cycle results from the high
bandwidth ($\omega_n=5$~rad/s) combined with the camera delay, the
clearance of the linkage joints and the dead zones. It motivated the
final tuning with $\omega_n=3$~rad/s (Section~\ref{sec:lqr}).

\begin{figure}[H]
    \centering
    \includegraphics[width=0.72\textwidth]{FOTOS/Fig12_servo_commands_test1.png}
    \caption{Control effort in integration run 1 with the preliminary
    tuning: servo commands (top: servo X, bottom: servo Y, in degrees).
    After the transient, the servos keep oscillating with an amplitude of
    about $\pm20$--$30^\circ$.}
    \label{fig:servo}
\end{figure}

\subsection{System test on flat floor}
\label{sec:flat}
The robot followed a C-shaped path, with straight sections and changes
of direction, and returned to its initial position, while the platform
balanced the ball. The test was repeated three times
(Table~\ref{tab:flat}).

\begin{table}[H]
\centering
\caption{System test on flat floor along a C-shaped path (final tuning,
reference at the center).}
\label{tab:flat}
\small
\begin{tabular}{@{}lccccccc@{}}
\toprule
\textbf{Run} & \textbf{Duration} & \textbf{Mean $|e_x|$} & \textbf{Mean $|e_y|$} & \textbf{Mean dist.} & \textbf{Mean dist.} & \textbf{Max.\ dist.} & \textbf{Time within} \\
 & [s] & [cm] & [cm] & [cm] & [\% of plate] & [cm] & \textbf{2 cm} [\%] \\
\midrule
1 & 206 & 0.67 & 0.60 & 0.99 & 5.0 & 4.45 & 92 \\
2 & 509 & 0.66 & 0.67 & 1.04 & 5.2 & 6.07 & 91 \\
3 & 240 & 0.57 & 0.52 & 0.85 & 4.3 & 3.14 & 97 \\
\midrule
Mean $\pm$ s.d. & -- & $0.63\pm0.06$ & $0.60\pm0.08$ & $0.96\pm0.10$ & $4.8\pm0.5$ & $4.55\pm1.47$ & $93\pm3$ \\
\bottomrule
\end{tabular}
\end{table}

The mean distance to the center was 0.96~cm (4.8\,\% of the plate side)
and the ball stayed within 2~cm of the center 93\,\% of the time. The
errors on X and Y are similar (0.63 and 0.60~cm), because on flat floor
there is no dominant disturbance in one direction. The distance exceeded
2~cm only in short peaks associated with the starts and turns of the
robot, after which the ball returned below 2~cm in one to two seconds.
The largest value, 6.07~cm in run 2, occurred in the first seconds of
the run, and the ball returned to the center without falling. When the
robot stopped at the end of a run, the distance remained at about
0.3~cm, inside the 0.4~cm dead zone of the controller.
Fig.~\ref{fig:flatpos} shows the position of the ball in the three
runs, and Fig.~\ref{fig:xy}a the corresponding trajectories on the
plate: the ball stays in a region of about $\pm2$~cm around the center,
with similar spread in X and Y.

\begin{figure}[H]
    \centering
    \includegraphics[width=\textwidth]{FOTOS/Fig14_flat_position.png}
    \caption{Ball position relative to the reference, $e_x$ (blue) and
    $e_y$ (orange), in the three runs on flat floor.}
    \label{fig:flatpos}
\end{figure}

\subsection{System test on the 10° ramp}
\label{sec:ramp}
In each of the four runs the robot climbed and descended the
$10^\circ$ ramp, and then moved back over it to its initial position
(videos in Table~\ref{tab:files_own}). The LabVIEW front panel showed a
forward displacement of about 3~m and a return to the origin with a
final position error of about 0.05~m, read from the screen recording.
Table~\ref{tab:ramp} shows the results for the ball.

\begin{table}[H]
\centering
\caption{System test on the $10^\circ$ ramp: climb, descent and reverse
motion (final tuning, reference at the center).}
\label{tab:ramp}
\small
\begin{tabular}{@{}lccccccc@{}}
\toprule
\textbf{Run} & \textbf{Duration} & \textbf{Mean $|e_x|$} & \textbf{Mean $|e_y|$} & \textbf{Mean dist.} & \textbf{Mean dist.} & \textbf{Max.\ dist.} & \textbf{Time within} \\
 & [s] & [cm] & [cm] & [cm] & [\% of plate] & [cm] & \textbf{2 cm} [\%] \\
\midrule
1 & 101 & 0.48 & 0.94 & 1.15 & 5.8 & 3.94 & 82 \\
2 & 97 & 0.40 & 0.89 & 1.05 & 5.3 & 7.00 & 81 \\
3 & 108 & 0.84 & 1.00 & 1.42 & 7.1 & 4.72 & 76 \\
4 & 73 & 0.67 & 1.23 & 1.54 & 7.7 & 4.81 & 71 \\
\midrule
Mean $\pm$ s.d. & -- & $0.60\pm0.20$ & $1.01\pm0.15$ & $1.29\pm0.23$ & $6.5\pm1.1$ & $5.12\pm1.32$ & $78\pm5$ \\
\bottomrule
\end{tabular}
\end{table}

On the ramp the mean distance increased to 1.29~cm (6.5\,\% of the
plate side), only 0.33~cm more than on flat floor, and the time within
2~cm decreased from 93 to 78\,\%. The error on Y, which is the direction
of motion of the robot, was almost twice the error on X (1.01 against
0.60~cm), whereas on flat floor both were similar. The slope is
therefore the dominant disturbance, together with the accelerations at
the entry and exit of each inclined section. This is visible in
Fig.~\ref{fig:xy}b, where the trajectories on the ramp are elongated in
the Y direction, and in Fig.~\ref{fig:ramppos}: in each run, $e_y$
shows plateaus of 2--3~cm (up to 4--5~cm in runs 2 and 4) that last
5--10~s, while the robot is on an inclined section. Between them, the
ball returns close to the center. The maximum distance, 7.0~cm in run
2, occurred in one of these transients, and the ball stayed on the
plate.

\begin{figure}[H]
    \centering
    \includegraphics[width=\textwidth]{FOTOS/Fig15_ramp_position.png}
    \caption{Ball position relative to the reference, $e_x$ (blue) and
    $e_y$ (orange), in the four runs on the $10^\circ$ ramp. The flat
    segments at the start of runs 2 and 4 correspond to the ball at
    rest inside the dead zone before the robot started to move.}
    \label{fig:ramppos}
\end{figure}

\begin{figure}[H]
    \centering
    \includegraphics[width=0.9\textwidth]{FOTOS/Fig17_xy_trajectories.png}
    \caption{Trajectory of the ball on the plate in (a) the three runs
    on flat floor and (b) the four runs on the $10^\circ$ ramp. Each color
    is one run.}
    \label{fig:xy}
\end{figure}

\subsection{Control effort}
\label{sec:effort}
The control effort was evaluated with the commanded plate tilt
$\theta_d$ (Eq.~\eqref{eq:law}) and with the
servo angles recorded in the logs (Table~\ref{tab:effort} and
Fig.~\ref{fig:effort}).

\begin{table}[H]
\centering
\caption{Control effort in the system tests (mean of the runs). RMS:
root mean square; saturation: percentage of samples with
$|\theta_d|\geq9.9^\circ$.}
\label{tab:effort}
\small
\begin{tabular}{@{}lcccccc@{}}
\toprule
\textbf{Test} & \textbf{RMS $\theta_{d,x}$} & \textbf{RMS $\theta_{d,y}$} & \textbf{Max.\ $|\theta_{d,y}|$} & \textbf{Saturation} & \textbf{Servo Y range used} & \textbf{Mean servo speed} \\
 & [deg] & [deg] & [deg] & [\%] & (1st--99th pct.) [deg] & [deg/s] \\
\midrule
Flat floor & 0.89 & 0.99 & 7.4 & 0.07 & 58--113 & 165 \\
$10^\circ$ ramp & 1.19 & 2.23 & 7.6 & 0.00 & 60--129 & 159 \\
\bottomrule
\end{tabular}
\end{table}

\begin{figure}[H]
    \centering
    \includegraphics[width=\textwidth]{FOTOS/Fig16_ramp_effort.png}
    \caption{Control effort on the $10^\circ$ ramp: commanded plate tilt
    about X (blue) and Y (orange) in the four runs.}
    \label{fig:effort}
\end{figure}

On the ramp, the RMS tilt about Y was 2.3 times that on flat floor
(2.23 against $0.99^\circ$), while about X it increased only slightly,
which agrees with the larger error on Y. Fig.~\ref{fig:effort} shows
that, on each inclined section, $\theta_{d,y}$ forms a plateau of
$4$--$6^\circ$, with the sign of the slope, and returns to about zero on
flat sections. The plateaus coincide in time with the offsets of
Fig.~\ref{fig:ramppos}. The data show that neither the controller nor
the mechanism saturated: $|\theta_d|$ never reached the $10^\circ$
limit, and the servo Y used 60--$129^\circ$, about $-37$ to $+32^\circ$
from its leveled position, within the $\pm50^\circ$ software range. The
offset is explained instead by the limit of $3^\circ$ on the integral
term (Section~\ref{sec:lqr}): with the integral at its limit, the
proportional term must provide the rest of the tilt, which requires a
steady error of
\begin{equation}
e_{ss}\approx\frac{\theta_{req}-\theta_{I,max}}{k_1}+e_{dz}
\approx\frac{4.6-3}{1.12}+0.4\approx1.8~\mathrm{cm},
\label{eq:offset}
\end{equation}
where $\theta_{req}\approx4.6^\circ$ is the mean commanded tilt measured
on the slope (run 1), $\theta_{I,max}=3^\circ$ the limit of the integral
term, $k_1$ the gain of Eq.~\eqref{eq:final} and $e_{dz}=0.4$~cm the
dead zone. This estimate is close to the measured offsets of
2--3~cm. Raising the integral limit to about $6^\circ$ would cancel
this offset; the limit was kept low to avoid windup during the starts
and turns of the robot. The mean servo speed was similar in both
conditions (165 and 159~deg/s), so the slope increases the amplitude of
the tilt but not the activity of the servos.

\subsection{Summary and comparison with the acceptance criteria}
Fig.~\ref{fig:summary} compares all the runs and Table~\ref{tab:accept}
summarizes them against the acceptance criteria.

\begin{figure}[H]
    \centering
    \includegraphics[width=\textwidth]{FOTOS/Fig13_summary_runs.png}
    \caption{Results of all the runs. (a) Mean distance between the ball
    and the reference; the dashed line is the 2~cm acceptance limit.
    (b) Percentage of time with the ball within 2~cm of the reference in
    the system tests.}
    \label{fig:summary}
\end{figure}

\begin{table}[H]
\centering
\caption{Comparison with the acceptance criteria (mean values per test).}
\label{tab:accept}
\small
\begin{tabularx}{\textwidth}{@{}L{3.6cm}YYYL{2.4cm}@{}}
\toprule
\textbf{Test} & \textbf{Ball fell} & \textbf{Mean distance [cm] (\% of plate)} & \textbf{Time within 2~cm} & \textbf{Criteria met} \\
\midrule
Integration (bench) & 0 of 5 runs & 0.56 (2.8\,\%) & -- & Yes \\
System, flat floor & 0 of 3 runs & 0.96 (4.8\,\%) & 93\,\% & Yes \\
System, $10^\circ$ ramp & 0 of 4 runs & 1.29 (6.5\,\%) & 78\,\% & Yes \\
\bottomrule
\end{tabularx}
\end{table}

\subsection{Capabilities and limitations}
\label{sec:limits}
\textbf{Capabilities.}
\begin{itemize}
    \item The ball was kept on the plate in all 12 runs, with a mean
    distance to the reference below 2~cm in every run, both with the base
    stopped and with the robot moving on flat floor and on a $10^\circ$
    ramp.
    \item The balancing loop works independently of the mobile base and
    uses only the camera, so it is not affected by the slope or by the
    accelerations of the robot as an IMU-based loop would be.
    \item The reference can be changed at run time to any point at least
    2~cm from the edges of the plate.
\end{itemize}
\textbf{Limitations and proposed improvements.} Table~\ref{tab:limitations}
relates each limitation to its measured effect and to a proposed
improvement.

\begin{table}[H]
\centering
\caption{Limitations of the prototype, observed effect and proposed improvement.}
\label{tab:limitations}
\small
\begin{tabularx}{\textwidth}{@{}L{3.0cm}YY@{}}
\toprule
\textbf{Limitation} & \textbf{Observed effect} & \textbf{Proposed improvement} \\
\midrule
Integral limited to $3^\circ$ & Offset of 2--3~cm on the slope; mean distance 1.29~cm on the ramp vs 0.96~cm on flat floor (Eq.~\eqref{eq:offset}) & Raise the limit to about $6^\circ$ with an anti-windup that resets it during turns \\
Tilt range of the linkage & About $9.3^\circ$ (X) and $13.3^\circ$ (Y); servo Y used up to $37^\circ$ of its $\pm50^\circ$ on the $10^\circ$ ramp & Longer crank or rod attached closer to the center for steeper slopes, at the cost of resolution \\
Servo activity & Mean servo speed of about 160~deg/s even on flat floor; wear and consumption & Larger dead zone or a filter on the command near the reference \\
Encoder resolution & One count is 0.218~rad/s at 92~Hz; speed error of about 5\,\% for any tuning & Measure the period between pulses or average over a longer window \\
Odometry only & Real final error of 1.16~cm vs 0.01~m odometric tolerance; slip not corrected & External reference (floor camera or markers) \\
Color-based vision & Requires the orange ball and lighting within the tuned HSV range; accuracy of the fixed homography not measured, and its error grows towards the edges when the plate tilts & Automatic HSV calibration; measure the homography error with the ball at known positions \\
Enable line & GPIO27 has no pull-down; the start-up order must be followed & 10~k$\Omega$ pull-down resistor \\
Number of runs & 3 runs on flat floor and 4 on the ramp, one prototype and one ball & More runs, other balls and slopes to generalize the results \\
\bottomrule
\end{tabularx}
\end{table}

% =============================================================================
% END MATTER
% =============================================================================
\section{Conclusions}
This article presented an open, reproducible ball-and-plate platform
mounted on a differential-drive robot, built from TETRIX parts, three
3D-printed parts and commercial electronics, with a total cost of
USD 1,398. The architecture separates the balancing loop (Raspberry
Pi 4B, camera, Kalman filter and LQR with integral action) from the
motion of the base (NI myRIO with a PD wheel-speed loop), linked by a
digital enable line and a UART channel. Using only the camera in the
balancing loop made it insensitive to the slope and to the
accelerations of the robot. In the system tests the ball never fell;
the mean distance to the reference was $0.96\pm0.10$~cm on flat floor
and $1.29\pm0.23$~cm on a $10^\circ$ ramp, with 93 and 78\,\% of the
time within 2~cm, and the base returned to its start with a real error
of $1.16\pm0.27$~cm. The main limitations are the offset on the slope
caused by the limit of the integral term, the encoder resolution and
the small number of runs. The platform is suited for teaching and
experimentation in state estimation, optimal control and the
integration of distributed embedded systems.

\section*{Ethics statements}
This work did not involve human subjects, animal experiments or data
collected from social media platforms.

\section*{CRediT author statement}
\textbf{Angel Garcia:} Software (camera interface, ball detection and
ball-and-plate control on the Raspberry Pi), Methodology, Investigation
(wiring and connections of the Raspberry Pi and the ball-and-plate
subsystem), Validation, Data curation, Writing -- review and editing.
\textbf{Thomas Flores:} Methodology (mechanical design), Visualization
(CAD models), Writing -- review and editing.
\textbf{Alejandro Guerra:} Methodology (mechanical design),
Visualization (CAD models), Writing -- review and editing.
\textbf{Carlos Hilario:} Resources, Investigation (mechanical
construction of the robot), Writing -- review and editing.
\textbf{Alex Valenzuela:} Software (motion control of the mobile base in
LabVIEW), Investigation (wiring and connections of the mobile base),
Validation, Writing -- review and editing.

\section*{Acknowledgements}
The authors thank the professors of the Mechatronic Systems Design course
and the laboratory staff of UTEC for their guidance and for access to the
laboratory equipment.

\section*{Use of AI tools}
Claude (Anthropic) was used to help write and edit the text of this
manuscript, format the \LaTeX{} source, suggest candidate literature
references, and generate the annotated versions of the CAD figures and
the diagrams (Figs.~\ref{fig:vmodel}, \ref{fig:proto_cad},
\ref{fig:views}, \ref{fig:mech}, \ref{fig:wiring}, \ref{fig:control},
\ref{fig:vision}, \ref{fig:color} and \ref{fig:interface}) from material
provided by the authors, to propose the configuration of the UART link
and the enable line, and to write the Python scripts that segment the
logged data and compute the metrics and plots of
Section~\ref{sec:validation}, including the comparison of the
wheel-speed tunings. The scores of the Pugh matrix were
reviewed by the authors. The
hardware, CAD models, control software, experiments and experimental
data were produced by the authors. The authors reviewed the text and
figures against their prototype, code and datasheets, verified the
references in the original sources, and are responsible for the final
content.

% =============================================================================
% REFERENCES
% =============================================================================
\begin{thebibliography}{99}

\bibitem{Awtar2002}
S. Awtar, C. Bernard, N. Boklund, A. Master, D. Ueda, and K. Craig,
``Mechatronic design of a ball-on-plate balancing system,''
\textit{Mechatronics}, vol. 12, no. 2, pp. 217--228, 2002,
doi: 10.1016/S0957-4158(01)00062-9.

\bibitem{ChengTsai2016}
C.-C. Cheng and C.-H. Tsai,
``Visual servo control for balancing a ball-plate system,''
\textit{International Journal of Mechanical Engineering and Robotics Research},
vol. 5, no. 1, pp. 28--32, 2016, doi: 10.18178/ijmerr.5.1.28-32.

\bibitem{Park2003}
J. H. Park and Y. J. Lee,
``Robust visual servoing for motion control of the ball on a plate,''
\textit{Mechatronics}, vol. 13, no. 7, pp. 723--738, 2003.

\bibitem{Hauser1992}
J. Hauser, S. Sastry, and P. Kokotovi\'c,
``Nonlinear control via approximate input-output linearization: the ball
and beam example,''
\textit{IEEE Transactions on Automatic Control}, vol. 37, no. 3,
pp. 392--398, 1992.

\bibitem{Heins2023}
A. Heins and A. P. Schoellig,
``Keep it upright: Model predictive control for nonprehensile object
transportation with obstacle avoidance on a mobile manipulator,''
\textit{IEEE Robotics and Automation Letters}, vol. 8, no. 12,
pp. 7986--7993, 2023, doi: 10.1109/LRA.2023.3324520.

\bibitem{Selvaggio2023}
M. Selvaggio, A. Garg, F. Ruggiero, G. Oriolo, and B. Siciliano,
``Non-prehensile object transportation via model predictive
non-sliding manipulation control,''
\textit{IEEE Transactions on Control Systems Technology}, 2023.

\bibitem{Brei2024}
Z. Brei, J. Michaux, B. Zhang, P. Holmes, and R. Vasudevan,
``Serving time: Real-time, safe motion planning and control for
manipulation of unsecured objects,''
\textit{IEEE Robotics and Automation Letters}, vol. 9, no. 3,
pp. 2383--2390, 2024.

\bibitem{Kaheman2023}
K. Kaheman, U. Fasel, J. J. Bramburger, B. Strom, J. N. Kutz, and
S. L. Brunton,
``The experimental multi-arm pendulum on a cart: A benchmark system for
chaos, learning, and control,''
\textit{HardwareX}, vol. 15, p. e00465, 2023,
doi: 10.1016/j.ohx.2023.e00465.

\bibitem{Myers2020}
A. D. Myers, J. R. Tempelman, D. Petrushenko, and F. A. Khasawneh,
``Low-cost double pendulum for high-quality data collection with
open-source video tracking and analysis,''
\textit{HardwareX}, vol. 8, p. e00138, 2020.

\bibitem{VDI2206}
VDI 2206, \textit{Design Methodology for Mechatronic Systems}.
Berlin, Germany: Beuth Verlag, 2004.

\bibitem{Pugh1991}
S. Pugh,
\textit{Total Design: Integrated Methods for Successful Product Engineering}.
Wokingham, UK: Addison-Wesley, 1991.

\bibitem{Suzuki1985}
S. Suzuki and K. Abe,
``Topological structural analysis of digitized binary images by border
following,'' \textit{Computer Vision, Graphics, and Image Processing},
vol. 30, no. 1, pp. 32--46, 1985.

\bibitem{Siegwart2011}
R. Siegwart, I. R. Nourbakhsh, and D. Scaramuzza,
\textit{Introduction to Autonomous Mobile Robots}, 2nd ed.
Cambridge, MA, USA: MIT Press, 2011.

\bibitem{Anderson1990}
B. D. O. Anderson and J. B. Moore,
\textit{Optimal Control: Linear Quadratic Methods}.
Englewood Cliffs, NJ, USA: Prentice Hall, 1990.

\bibitem{Astrom2021}
K. J. \AA str\"om and R. M. Murray,
\textit{Feedback Systems: An Introduction for Scientists and Engineers},
2nd ed. Princeton, NJ, USA: Princeton University Press, 2021.

\bibitem{Kalman1960}
R. E. Kalman,
``A new approach to linear filtering and prediction problems,''
\textit{Journal of Basic Engineering}, vol. 82, no. 1, pp. 35--45, 1960.

\bibitem{BarShalom2001}
Y. Bar-Shalom, X. R. Li, and T. Kirubarajan,
\textit{Estimation with Applications to Tracking and Navigation}.
New York, NY, USA: Wiley, 2001.

\bibitem{Smith1957}
O. J. M. Smith,
``Closer control of loops with dead time,''
\textit{Chemical Engineering Progress}, vol. 53, no. 5, pp. 217--219, 1957.

\bibitem{Gonzalez2018}
R. C. Gonzalez and R. E. Woods,
\textit{Digital Image Processing}, 4th ed. New York, NY, USA: Pearson, 2018.

\bibitem{Hartley2004}
R. Hartley and A. Zisserman,
\textit{Multiple View Geometry in Computer Vision}, 2nd ed.
Cambridge, UK: Cambridge University Press, 2004.

\bibitem{OpenCV}
G. Bradski,
``The OpenCV library,''
\textit{Dr. Dobb's Journal of Software Tools}, vol. 25, no. 11,
pp. 120--125, 2000.

\bibitem{Picamera2}
Raspberry Pi Ltd.,
``The Picamera2 Library,'' manual. [Online]. Available:
\url{https://datasheets.raspberrypi.com/camera/picamera2-manual.pdf}

\bibitem{RPiCamera}
Raspberry Pi Ltd.,
``Camera -- Raspberry Pi Documentation.'' [Online]. Available:
\url{https://www.raspberrypi.com/documentation/accessories/camera.html}

\bibitem{RPi4}
Raspberry Pi Ltd.,
``Raspberry Pi 4 Model B,'' product brief. [Online]. Available:
\url{https://datasheets.raspberrypi.com/rpi4/raspberry-pi-4-product-brief.pdf}

\bibitem{RPiUART}
Raspberry Pi Ltd.,
``Configuration -- UARTs,'' Raspberry Pi Documentation. [Online]. Available:
\url{https://www.raspberrypi.com/documentation/computers/configuration.html}

\bibitem{OV5647}
OmniVision Technologies, Inc.,
``OV5647: 1/4'' color CMOS QSXGA (5 megapixel) image sensor with
OmniBSI technology,'' preliminary specification, version 1.0, 2009.

\bibitem{MPU6500}
InvenSense Inc.,
``MPU-6500 Product Specification,'' document PS-MPU-6500A-01.

\bibitem{Pedley2013}
M. Pedley,
``Tilt sensing using a three-axis accelerometer,''
Freescale Semiconductor, Application Note AN3461, rev. 6, 2013.

\bibitem{PCA9685}
NXP Semiconductors,
``PCA9685: 16-channel, 12-bit PWM Fm+ I$^2$C-bus LED controller,''
product data sheet, rev. 4, 2015.

\bibitem{MG996R}
Tower Pro,
``MG996R High Torque Metal Gear Dual Ball Bearing Servo,'' datasheet.

\bibitem{XL4015}
XLSEMI,
``XL4015: 180~kHz 36~V 5~A switching current buck DC to DC converter,''
datasheet.

\bibitem{XL4016}
XLSEMI,
``XL4016: 180~kHz 40~V 8~A switching current buck DC to DC converter,''
datasheet.

\bibitem{myRIO}
National Instruments,
``User Guide and Specifications: NI myRIO-1900.''

\bibitem{MDD10A}
Cytron Technologies,
``MDD10A Rev2.0: 10~A dual channel DC motor driver,'' user's manual.

\end{thebibliography}

% =============================================================================
% APPENDIX
% =============================================================================
\section*{Appendix A. Additional figures}
\addcontentsline{toc}{section}{Appendix A. Additional figures}
\setcounter{figure}{0}
\renewcommand{\thefigure}{A.\arabic{figure}}

\begin{figure}[H]
    \centering
    \includegraphics[width=\textwidth]{FOTOS/FigA1_base_all_tunings.png}
    \caption{Right-wheel speed (blue), reference (dashed) and PWM (orange)
    during the first 20~s of each wheel-speed test
    (Section~\ref{sec:wheel}).}
    \label{fig:appA1}
\end{figure}

\begin{figure}[H]
    \centering
    \includegraphics[width=\textwidth]{FOTOS/FigB1_flat_error.png}
    \caption{Distance between the ball and the reference in the flat-floor runs
    (Section~\ref{sec:flat}).}
    \label{fig:appA2}
\end{figure}

\begin{figure}[H]
    \centering
    \includegraphics[width=\textwidth]{FOTOS/FigB2_ramp_error.png}
    \caption{Distance between the ball and the reference in the ramp runs
    (Section~\ref{sec:ramp}).}
    \label{fig:appA3}
\end{figure}

\begin{figure}[H]
    \centering
    \includegraphics[width=\textwidth]{FOTOS/FigB3_flat_effort.png}
    \caption{Commanded plate tilt (control effort of the LQR) in the
    flat-floor runs, for comparison with the ramp in
    Fig.~\ref{fig:effort}.}
    \label{fig:appA4}
\end{figure}
\clearpage
\setcounter{figure}{0}
\renewcommand{\thefigure}{\arabic{figure}}

% =============================================================================
% BIOGRAPHIES
% =============================================================================
\section*{Biographies}

Although each author had different roles (see the CRediT statement), the workload was shared equally (20\,\% each).

\authorbio{Garcia Motta, Angel Giovani}{FOTOS/Angel.jpg}{%
is an undergraduate student in Mechatronics Engineering at Universidad de
Ingenier\'ia y Tecnolog\'ia (UTEC), Lima, Peru. His interests include
robotics, automation, control systems and mechatronic design. In this
project he contributed mainly to the camera interface, the ball detection and the control of the ball-and-plate platform on the Raspberry Pi, and the wiring of the ball-and-plate electronics.
}

\authorbio{Flores Baltodano, Thomas Mat\'ias}{FOTOS/Thomas.jpg}{%
is an undergraduate student in Mechatronics Engineering at Universidad de
Ingenier\'ia y Tecnolog\'ia (UTEC), Lima, Peru. In this project he
contributed mainly to the CAD design of the robot and of the ball-and-plate mechanism.
}

\authorbio{Guerra Figueroa, Alejandro}{FOTOS/Alejandro.jpeg}{%
is an undergraduate student in Mechatronics Engineering at Universidad de
Ingenier\'ia y Tecnolog\'ia (UTEC), Lima, Peru. In this project he
contributed mainly to the CAD design of the robot and of the ball-and-plate mechanism.
}

\authorbio{Hilario Mayta, Carlos Eduardo}{FOTOS/author4.png}{%
is an undergraduate student in Mechatronics Engineering at Universidad de
Ingenier\'ia y Tecnolog\'ia (UTEC), Lima, Peru. In this project he
contributed mainly to the mechanical construction of the robot.
}

\authorbio{Valenzuela Quispe, Alex Claudio}{FOTOS/alex_foto.jpeg}{%
is an undergraduate student in Mechatronics Engineering at Universidad de
Ingenier\'ia y Tecnolog\'ia (UTEC), Lima, Peru. In this project he
contributed mainly to the motion control of the mobile base and its wiring and connections.
}

\end{document}
